%!TEX root = ../main.tex
% =============================================================================
%  附录（格式规范第四条：页数不限，必须打印并与正文装订在一起提交）
%  附录必备八项（依据 format2026.doc 第五条、第十一条）：
%    1 支撑材料文件清单        -> 附录 A  ← 必需
%    2 全部完整、可运行的源程序 -> 附录 E  ← 缺少可能被取消评奖资格
%    3 运行环境与依赖版本      -> 附录 B
%    4 运行说明               -> 附录 B
%    5 随机种子清单            -> 附录 H
%    6 结果文件与论文表格对应表 -> 附录 C
%    7 数据字典               -> 附录 F
%    8 较长的推导              -> 附录 G
%  ⚠️ 若确实没有用到程序，须写明「本论文没有用到程序」（第五条）
%  ⚠️ 若确实没有支撑材料，须注明「本论文没有支撑材料」（第十一条）
%  ⚠️ 附录中同样不得出现任何身份信息（第六条）：校名、队伍编号、姓名、所在区域一律不写
%  ⚠️ 附录 A 的文件清单与附录 E 的 \codefile 清单都按工程**实际文件**逐项写，
%     新增脚本或结果文件必须同步补入这两处（附录 E 缺文件会被判为红线）。
% =============================================================================

\section{支撑材料文件清单}

支撑材料压缩包 \texttt{support.zip} 打包本工程「可复现必需」的文件：论文源文件、全部源程序、
五个结果文件、八张数据图的矢量文件，以及作为论文数字来源的汇总类中间结果。
下表与本工程下的实际文件逐项对应；附录~\ref{sec:app-code} 所列源程序与本表的 \texttt{code/} 完全一致。

\begin{lstlisting}[style=plain,basicstyle=\ttfamily\footnotesize\songti]
support.zip
+-- main.tex                        论文装配（宏包输入、附录起点、正文结束标记）
+-- preamble.tex                    导言区：格式合规配置与写作宏（口径框/证据号/占位符）
+-- build.sh                        编译入口（自动探测 latexmk+xelatex / xelatex / tectonic）
+-- sections/                       正文与附录源文件，13 个
|   +-- 00_abstract.tex             摘要专用页
|   +-- 01_problem.tex              一、问题重述与分析
|   +-- 02_assumptions.tex          二、模型假设与符号说明
|   +-- 03_data.tex                 三、数据预处理与特征分析
|   +-- 04_q1.tex                   四、问题 1 模型
|   +-- 05_q2.tex                   五、问题 2 模型
|   +-- 06_q3.tex                   六、问题 3 模型
|   +-- 07_q4.tex                   七、问题 4 模型
|   +-- 08_results.tex              八、结果与分析
|   +-- 09_validation.tex           九、模型检验与灵敏度
|   +-- 10_evaluation.tex           十、模型评价与推广
|   +-- refs.tex                    参考文献（GB/T 7714）
|   +-- 90_appendix.tex             本附录
|
+-- tools/
|   +-- check_paper.py              提交前合规自查（--strict 为定稿模式）
|   +-- check_numbers.py            论文数字与数值台账的机械对拍
|   +-- verify_ref.py               对照稿红队审查的只读反算脚本（复现 notes/F 的全部断言）
|   +-- package_support.sh          本压缩包的打包脚本（保证本清单与实际文件一致）
|
+-- code/                           全部源程序，与附录 E 逐项一致
|   +-- config.py                   口径、参数与附件读取（全文唯一参数来源）
|   +-- lp_engine.py                单日两阶段 LP、滚动子问题 LP、场景削减
|   +-- q1.py  q2.py  q3.py  q4.py  四问求解与结果文件写出
|   +-- results_io.py               结果工作表写出器（按附件 5 模板的列序）
|   +-- exp_scan.py                 容量/功率/效率扫描与报童经验覆盖率
|   +-- exp_sobol.py                全局灵敏度（Sobol 指数）
|   +-- exp_na.py                   非预期性置换检验
|   +-- exp_oos.py                  样本外回测（训练窗构造情景集、测试窗回放）
|   +-- exp_benchmark.py            全年单一 LP 基准实验（完全信息参照）
|   +-- audit.py                    独立审计 A1--A10（不共用建模费用函数）
|   +-- paper_numbers.py            汇总数值台账 paper_numbers.json
|   +-- plot_figures.py             八张数据图的绘图脚本
|   +-- plot_figures_v2.py          v2 新增五张数据图（相依结构/生成器对照/风险画像/鲁棒前沿/分解）
|   +-- copula.py                  二元 Copula 工具库与 C 藤 pair-Copula（密度/h 函数/拟合/抽样）
|   +-- scen_copula.py             KDE + Copula 联合情景生成器与保真度度量
|   +-- exp_copula.py              相依结构拟合、四种情景生成器样本外对照、分时段 CVaR 风险画像
|   +-- exp_chance.py              机会约束对照（安全边际分位数扫描与报童互证）
|   +-- exp_dro.py                 盒式分布鲁棒对照（保守度—费用曲线与最坏情形评估）
|   +-- exp_loaderr.py             假设 A9 的量化：负载预报误差 0%/5%/10% 三档重跑问题 2
|   +-- decomp.py                  L 形（Benders）分解对照（解析多割、与直接求解对拍）
|   +-- run_all.py                  一键复现
|   +-- requirements.txt            依赖清单（版本锁定）
|
+-- results/                        五个结果文件（工作表名与列序按附件 5 模板）
|   +-- result1.xlsx                问题 1（典型日）：计划购电量、充放电量
|   +-- result2.xlsx                问题 2：计划购电量、充放电量、紧急购电量
|   +-- result3.xlsx                问题 3：计划购电量、调整购电量、充放电量、紧急购电量
|   +-- result4-2.xlsx              问题 4-2：与 result2 同结构
|   +-- result4-3.xlsx              问题 4-3：与 result3 同结构
|
+-- figures/                        图 1 至图 10
|   +-- fig01_timeline.tex          图 1 决策时序图（TikZ 源文件）
|   +-- fig02_dependency.tex        图 2 问题依赖图（TikZ 源文件）
|   +-- fig03_q1_dual_price.pdf     电价、水价值与阈值型策略（问题 1）
|   +-- fig04_q1_capacity_power.pdf 容量与充放功率上限双向扫描（问题 1）
|   +-- fig05_cost_waterfall.pdf    问题 1 成本分解瀑布
|   +-- fig06_voi.pdf               信息价值曲线（问题 3）
|   +-- fig07_q2_emergency.pdf      逐日紧急购电与报童两界（问题 2）
|   +-- fig08_q4_price.pdf          电价过程骨架与残差结构（问题 4）
|   +-- fig09_cvar_frontier.pdf     CVaR 风险--费用前沿与真实回测（问题 4）
|   +-- fig10_strategy_mechanism.pdf 最优策略机制（问题 1）
|   +-- README.md                   图规格、配色语义与出图约定
|
+-- intermediate/                   中间结果（论文数字的来源），分两组
|   +-- 随包提交（汇总与审计类，单个不超过 0.21 MB）
|   |   +-- audit.json  audit_report.md          审计 A1--A10（附录 D 的机器可读与可读版）
|   |   +-- paper_numbers.json                   数值台账（对拍基准，含四问、对偶、消融、Sobol、CVaR）
|   |   +-- q1.json  q1_dual.json  q1_dual2.json 问题 1 解与最优性证书
|   |   +-- q2_margin.json  exp_benchmark.json  安全边际分解、全年单一 LP 基准
|   |   +-- q4_price.json  q4_cvar.json          电价过程检验与情景保真度、CVaR 前沿
|   |   +-- exp_scan.json  exp_sobol.json  exp_na.json  扫描、灵敏度、非预期性
|   |   +-- full_year_lp.txt                     全年单一 LP（完全信息参照）的目标值
|   +-- 不随包（逐日解转储，单个 2.7--4.4 MB，可由 run_all.py 一键重建）
|       +-- q2_main.json  q2_noMargin.json  q2_cmp_R1.json  q2_cmp_R2.json  q2_cmp_R3.json
|       +-- q2_K120.json  q2_full_strat.json  q2_term_R1_eq.json  q2_term_R1_free.json
|       +-- q2_term_R3_free.json  q3_n1.json  q3_n2.json  q3_n4.json  q3_perf.json
|       +-- q4_2.json  q4_3.json
|
+-- notes/                          口径与台账文档，7 个
    +-- 写作交底.md                 本队实际模型口径（唯一口径说明）
    +-- 数值台账.md                 论文数字的表格版台账
    +-- A_坑点与验证.md  B_模型与算法.md  C_写作规范.md
    +-- D_行业与引用.md  E_知识点与模板.md
\end{lstlisting}

三点说明：

\begin{enumerate}
  \item \textbf{体积}：上述随包内容合计约 3.4~MB（其中五个结果文件约 2.1~MB、
        八张数据图约 0.3~MB），低于支撑材料压缩包 20~MB 的上限（格式规范第十一条）。
  \item \textbf{中间结果分级打包}：\texttt{intermediate/} 下第二组\textbf{逐日解转储}
        （\texttt{q2\_main.json}、\texttt{q2\_noMargin.json}、\texttt{q3\_n4.json}、\texttt{q4\_2.json} 等，
        单个 2.7--4.4~MB）可由 \texttt{python3 code/run\_all.py} 一键重建，
        且其中每个进入论文的数字都已汇总进 \texttt{paper\_numbers.json}，故不随压缩包提交；
        包内保留体积小、且为论文数字与审计直接依据的汇总类文件（上表已逐一列出）。
  \item \textbf{不随包提交的中间产物}：LaTeX 编译产物、Python 字节码缓存、以下划线开头的临时脚本与
        \texttt{.npy} 临时数组，均为可再生的工作文件，不属于支撑材料。
\end{enumerate}

\section{运行环境与运行说明}

\paragraph{运行环境}

\begin{table}[htbp]
  \centering
  \caption{本文的运行环境与依赖版本清单（对象：操作系统、解释器、数值与绘图库、求解器接口；
            范围：全部求解、审计与出图的实际执行环境；单位与量纲不适用；样本量为 1 套环境配置）}
  \label{tab:env}
  \footnotesize
  \begin{tabularx}{\linewidth}{L{3.4cm} L{4.6cm} Y}
    \toprule
    项目 & 版本 & 说明 \\
    \midrule
    操作系统 & macOS 15.7.5（arm64，Build 24G624） & 全部求解、审计与出图单机完成，未使用集群 \\
    Python & 3.13.12 & 唯一解释器 \\
    NumPy / pandas & 2.4.3 / 3.0.1 & 数组运算与附件 \texttt{.xlsx} 读写 \\
    SciPy & 1.18.1 & 线性规划 \texttt{linprog} 与统计检验 \\
    求解器 & SciPy \texttt{linprog(method=“highs”)}，HiGHS 后端（随 SciPy 1.18.1 分发）
      & 四问模型在场景离散后均为 LP，由 HiGHS 精确求解，未使用商业求解器 \\
    Matplotlib & 3.11.1 & 图 3--图 10 的数据图；图 1、图 2 为 TikZ 源文件，由 XeLaTeX 排版 \\
    scikit-learn & 1.9.0 & 仅用于 $k$-means 场景削减的对照路径 \\
    openpyxl & 3.1.5 & 结果文件 \texttt{.xlsx} 的读写 \\
    \bottomrule
  \end{tabularx}
\end{table}

依赖清单锁定在 \texttt{code/requirements.txt}（\texttt{numpy==2.4.3}、\texttt{pandas==3.0.1}、
\texttt{scipy==1.18.1}、\texttt{openpyxl==3.1.5}、\texttt{matplotlib==3.11.1}、
\texttt{scikit-learn==1.9.0}），版本号写确切值、不写区间。

\paragraph{运行说明}

\begin{enumerate}
  \item \textbf{依赖安装}：\texttt{python3 -m pip install -r code/requirements.txt}；
  \item \textbf{数据就位}：把附件 1--4 与附件 5 模板放在工作区的
        \texttt{CUMCM2026Problems(1)/C题/附件/}（模板置于其 \texttt{附件5/} 子文件夹）。
        该路径由 \texttt{code/config.py} 的 \texttt{ATTACH\_DIR} 给出，并由脚本自身位置反推，
        整个工程文件夹移动后无需改动代码；
  \item \textbf{依次运行四问}（在 \texttt{code/} 下执行，命令参数即论文主口径）：

\begin{lstlisting}[style=plain,basicstyle=\ttfamily\footnotesize\songti]
python3 q1.py                                                      # 问题 1：典型日 LP
python3 q2.py --tag main --K 60 --scen strat --info R3 --write    # 问题 2：334 天
python3 q3.py --epochs 0,6,12,18 --tag n4 --K 60 --write          # 问题 3：四次预报
python3 q4.py --mode q42 --K 64 --write                           # 问题 4-2：波动电价
python3 q4.py --mode q43 --K 64 --write                           # 问题 4-3：滚动
\end{lstlisting}

        \texttt{--write} 表示按附件 5 模板写出对应结果文件；
        条件项（论文中作为对照或检验给出）另有：\texttt{q4.py --mode tests --K 64}（电价过程检验与
        情景保真度）、\texttt{q4.py --mode cvar --K 64 --step 6}（CVaR 前沿）、
        \texttt{exp\_scan.py}（容量、功率与效率扫描、报童经验覆盖率）、
        \texttt{exp\_sobol.py}（Sobol 全局灵敏度）、\texttt{exp\_na.py}（非预期性置换检验）、
        \texttt{exp\_benchmark.py}（全年单一 LP 的完全信息基准）、
        \texttt{plot\_figures.py}（八张数据图）；
  \item \textbf{独立审计}：\texttt{python3 audit.py}。脚本只读原始附件与五个结果文件，
        独立复算 A1--A10，写出 \texttt{intermediate/audit.json} 与 \texttt{audit\_report.md}；
  \item \textbf{数值台账}：\texttt{python3 paper\_numbers.py}，把全部中间结果汇总为
        \texttt{intermediate/paper\_numbers.json}（论文数字的对拍基准）；
  \item \textbf{对拍与自查}：\texttt{python3 ../tools/check\_numbers.py} 列出论文中找不到
        台账出处的数字；\texttt{python3 ../tools/check\_paper.py --strict} 做定稿模式合规自查
        （占位符、禁令词、身份信息、缺图与缺源码、页数与摘要页）。
\end{enumerate}

\noindent\textbf{一键复现}：\texttt{python3 run\_all.py} 按上述顺序自动执行问题 1--4、
扫描、审计与台账；加 \texttt{--fast} 参数则跳过信息结构、场景数、VoI 与 CVaR 等对照实验。

\paragraph{实测耗时}

\begin{table}[htbp]
  \centering
  \caption{单机实测耗时（macOS 15.7.5 / arm64 / Python 3.13.12，与正文同一份代码与数据）}
  \label{tab:runtime}
  \footnotesize
  \begin{tabularx}{\linewidth}{L{4.0cm} L{2.0cm} Y}
    \toprule
    步骤 & 实测耗时 & 规模与说明 \\
    \midrule
    问题 1（典型日，确定性 LP） & 不足 0.1~s & 单日 144 槽；两次实测 0.014~s 与 0.008~s（纯求解时间） \\
    问题 2（固定价，两阶段） & 约 30~s & 334 天逐日两阶段子问题（$K=60$ 个场景） \\
    问题 3（固定价，四次预报） & 约 72~s & 四个预报时刻 $\times$ 334 天的滚动子问题 \\
    问题 4-2（波动电价，两阶段） & 约 50~s & 334 天逐日两阶段子问题（$K=64$ 个联合情景） \\
    问题 4-3（波动电价，四次预报） & 约 101~s & 四个预报时刻 $\times$ 334 天的滚动子问题 \\
    独立审计（A1--A10） & 约 13~s & 两次实测 12.9~s 与 16.7~s；其中交付物复算 3.3~s、逐日解复算 10.1~s \\
    \midrule
    合计（四问主链 $+$ 审计） & 约 4.4~min & 单机完成，无需集群或并行 \\
    \bottomrule
  \end{tabularx}
\end{table}

\noindent\textbf{输出文件}：\texttt{results/} 下的五个结果文件（\texttt{result1.xlsx}、
\texttt{result2.xlsx}、\texttt{result3.xlsx}、\texttt{result4-2.xlsx}、\texttt{result4-3.xlsx}）；
\texttt{figures/} 下的十三张矢量图（图 3--图 10 与 v2 新增的图 11--图 15）与两张 TikZ 源文件（图 1、图 2）；
\texttt{intermediate/} 下的审计报告、数值台账与各问中间结果。

\section{结果文件、论文表格与生成函数}\label{sec:app-filemap}

五个结果文件均按附件 5 模板的工作表名、列名与列序写出：\texttt{result1.xlsx} 为竖表
（一个时段一行，对应表~\bref{tab:q1-buy} 与表~\bref{tab:q1-soc}）；其余四个文件为横表
（一天一行，147 列 $=$ 日期 $+$ 144 个时段 $+$ 全天购电量 $+$ 全天购电费）。

\begin{table}[htbp]
  \centering
  \caption{结果文件 $\leftrightarrow$ 论文表格 $\leftrightarrow$ 生成函数 对应表}
  \label{tab:filemap}
  \footnotesize
  \setlength{\tabcolsep}{4pt}
  \begin{tabularx}{\linewidth}{L{3.4cm} L{2.4cm} L{4.0cm} Y}
    \toprule
    结果文件（工作表） & 论文位置 & 生成函数 & 行数与中间结果 \\
    \midrule
    result1.xlsx（计划购电量） & 表~\bref{tab:q1-buy}
      & \texttt{q1.solve\_q1} $\rightarrow$ \texttt{q1.write\_result1}
      & 144 行（1 行 1 个 10~min 槽）；\texttt{q1.json} \\
    result1.xlsx（充放电量） & 表~\bref{tab:q1-soc}
      & \texttt{q1.write\_result1}
      & 6 行（4~h 块）$+$ $0{:}00$ 与 $24{:}00$ 储电量；\texttt{q1.json} \\
    result2.xlsx（计划购电量） & 表~\ref{tab:q2-buy}、表~\bref{tab:q2-t1}、表~\bref{tab:summary}
      & \texttt{q2.run\_q2} $\rightarrow$ \texttt{results\_io.write\_result}
      & 334 行 $\times$ 147 列；\texttt{q2\_main.json} \\
    result2.xlsx（充放电量） & 表~\ref{tab:q2-soc}、表~\bref{tab:q2-t2} & 同上
      & 2\,004 行（334 天 $\times$ 6 个 4~h 块） \\
    result2.xlsx（紧急购电量） & 表~\ref{tab:q2-emergency}、表~\bref{tab:emergency} & 同上
      & 1\,156 行（缺口时段合并，同一日期可多行） \\
    result3.xlsx（计划购电量、调整购电量） & 表~\bref{tab:q3-epochs}（决策时刻与生效槽位）、表~\bref{tab:key4}
      & \texttt{q3.run\_q3} $\rightarrow$ \texttt{results\_io.write\_result}
      & 各 334 行 $\times$ 147 列；\texttt{q3\_n4.json} \\
    result3.xlsx（充放电量、紧急购电量） & 表~\bref{tab:key4}、表~\bref{tab:summary} & 同上
      & 2\,004 行 / 1\,190 行 \\
    result4-2.xlsx（计划购电量、充放电量、紧急购电量） & 表~\bref{tab:q4-summary}、表~\bref{tab:key4}、表~\bref{tab:summary}
      & \texttt{q4.run\_q4\_2} $\rightarrow$ \texttt{results\_io.write\_result}
      & 334 $\times$ 147、2\,004、1\,193 行；\texttt{q4\_2.json} \\
    result4-3.xlsx（计划购电量、调整购电量、充放电量、紧急购电量） & 表~\bref{tab:q4-summary}、表~\bref{tab:key4}
      & \texttt{q4.run\_q4\_3} $\rightarrow$ \texttt{results\_io.write\_result}
      & 334 $\times$ 147（两次）、2\,004、1\,187 行；\texttt{q4\_3.json} \\
    \bottomrule
  \end{tabularx}
\end{table}

表中「生成函数」为模块内的真实函数名：\texttt{q2.run\_q2}、\texttt{q3.run\_q3}、
\texttt{q4.run\_q4\_2}、\texttt{q4.run\_q4\_3} 完成求解，由各脚本的 \texttt{main()} 装配后
统一调用 \texttt{results\_io.write\_result} 写出工作簿；问题 1 的竖表由 \texttt{q1.solve\_q1} 求解、
\texttt{q1.write\_result1} 写出。论文其余图表不来自结果文件，其来源为：最优性证书取自
\texttt{q1\_dual.json} 与 \texttt{q1\_dual2.json}（\texttt{q1.py} 的对偶校验），
CVaR 前沿取自 \texttt{q4\_cvar.json}（\texttt{q4.py --mode cvar}，底层函数 \texttt{q4.solve\_cvar\_day}），
全局灵敏度取自 \texttt{exp\_sobol.json}，非预期性置换检验取自 \texttt{exp\_na.json}。

\section{独立审计报告}

审计由\textbf{非建模实现者}独立完成：脚本 \texttt{code/audit.py} 只读原始附件与本队写出的五个结果文件，
\textbf{不共用建模代码的费用函数与约束构造函数}，以保证「自证循环」不会发生。
机器可读结果在 \texttt{intermediate/audit.json}，可读报告在 \texttt{intermediate/audit\_report.md}。

\paragraph{审计口径与阈值}

\begin{itemize}
  \item \textbf{残差}：模型内部解（10~min 槽级、未取整）阈值 $10^{-6}$~kWh；
        交付物（结果文件只给 4~h 块级充放电量与 $0{:}00$、$24{:}00$ 两个储电量）的电量按
        $10^{-4}$~kWh 取整后回读，块级递推残差的上界为
        $0.5\times10^{-4}+0.9\times6\times0.5\times10^{-4}+6\times0.5\times10^{-4}/0.9+0.5\times10^{-4}=7.0\times10^{-4}$~kWh，
        故交付物的残差阈值取 $10^{-3}$~kWh
        （贴合交付精度的判定，不放松内部解的口径）；
  \item \textbf{越界数与违反数}：储电量越界、充放电功率越界、充放互斥违反、负电量，全部要求等于 $0$；
  \item \textbf{费用复算相对误差}：$\le10^{-4}$（费用只用结果文件与附件逐时段重算，不调用建模费用函数）。
\end{itemize}

\paragraph{逐项结果（正文口径）}

下表为正文第九章原先给出的\textbf{双列口径}审计表：左列为 10~min 级模型内部解，右列为结果文件
回读后的块级复算，便于对照“内部解”与“交付物”两层判定的差别。

\begin{table}[htbp]
  \centering
  \caption{独立审计项、阈值与实测结果（内部解：问题 1、2、3 的 10~min 级解；交付物：result1、result2、
           result3、result4-2、result4-3 共五个结果文件；窗口 2025-02-01 至 12-31）}
  \label{tab:audit}
  \scriptsize
  \setlength{\tabcolsep}{3pt}
  \begin{tabularx}{\linewidth}{C{0.62cm} L{3.95cm} L{3.35cm} L{3.35cm} C{0.95cm}}
    \toprule
    项 & 检查内容与阈值 & 内部解实测 & 交付物实测 & 判定 \\
    \midrule
    A1 & 功率/电量平衡残差 \(\le10^{-6}\)~kWh
      & 问题 1 \(1.49\times10^{-13}\)；问题 2、3 \(3.41\times10^{-13}\)~kWh
      & 交付物无逐槽 \(L,PV\) 列，平衡只能在内部解上核验（汇总列对拍见 A10） & PASS \\
    A2 & 储电量递推残差 \(\le10^{-6}\)~kWh
      & 问题 1 \(9.09\times10^{-13}\)；问题 2、3 \(1.09\times10^{-11}\)~kWh
      & 问题 1 \(3.64\times10^{-12}\)；问题 2、3、4-2、4-3 为
        \(2.96/2.72/2.11/2.61\times10^{-4}\)~kWh，为写入取整所致 & PASS\textsuperscript{注} \\
    A3 & 储电量边界：越界数 \(=0\)（区间 \([1200,10\,800]\)~kWh） & 0 & 0 & PASS \\
    A4 & 充放功率边界：内部解按单槽上限 \(833.3333\)~kWh；交付物按块级上限 \(20\,000\)~kWh & 0 & 0 & PASS \\
    A5 & 充放互斥：\(\min(c_t,d_t)\le10^{-6}\)
      & 违反槽数 0（充放同时槽数 0）
      & 交付物只有 4~h 聚合量，块内可同时有充放电，故块级不可检验 & PASS \\
    A6 & 周期/终端：问题 1 要求 \(S_0-S_{144}=0\)；问题 2--4 要求跨日衔接且 \(S_T\ge S_0\)
      & 问题 1 周期偏差 \(0.000\)；问题 2、3 末值低于初值的天数为 0
      & 五个结果文件末值低于初值的天数均为 0 & PASS \\
    A7 & 非负性：所有电量与费用 \(\ge-10^{-6}\) & 负购电槽数 0 & 负购电量数 0 & PASS \\
    A8 & 费用复算（计划 \(+\) 紧急 \(+\) 调整相关）相对误差 \(\le10^{-4}\)
      & 问题 1、2、3 的相对误差均为 0
      & 交付物层只能核验计划购电费（最大相对偏差 \(4.32\times10^{-7}\)，问题 4-2） & PASS \\
    A9 & 日期与结构：334 天、工作表名/列名/列序与模板一致
      & 334 天
      & 五个文件均为 334 行 \(\times\)147 列，日期范围 2025-02-01 至 12-31；紧急购电行数
        1156/1190/1193/1187 & PASS \\
    A10 & 全天汇总列与逐时段求和一致：相对误差 \(\le10^{-4}\) & ---
      & 全天购电量最大偏差 \(0.0012\)~kWh（问题 4-3）；全天购电费最大相对偏差
        \(4.32\times10^{-7}\)（问题 4-2） & PASS \\
    \bottomrule
  \end{tabularx}
  \\[1ex]
  {\footnotesize 表注：交付物 A2 的 \(10^{-4}\)~kWh 量级残差可由写入精度完全解释——结果写入对 6 个块级
  充放电量与 \(0{:}00\)/\(24{:}00\) 两个储电量按 4 位小数取整，取整误差上界为
  \(6\times5\times10^{-5}+2\times5\times10^{-5}=4\times10^{-4}\)~kWh，与实测值一致；该残差不是物理
  越界（A3、A4 越界数均为 0）。}
\end{table}

\paragraph{逐项结果}

\begin{table}[htbp]
  \centering
  \caption{独立审计 A1--A10（同时承载正文的 L1--L5 验证链条）：阈值、实测值与判定
           （「内」为 10~min 级模型内部解，「交」为结果文件块级）}
  \label{tab:audit-appendix}
  \footnotesize
  \setlength{\tabcolsep}{4pt}
  \begin{tabularx}{\linewidth}{C{0.9cm} L{3.4cm} L{2.9cm} Y C{1.2cm}}
    \toprule
    项 & 检查内容 & 阈值 & 实测值 & 判定 \\
    \midrule
    A1 & 电量平衡残差 & $\max\lvert E_t+PV_t\Delta t+d_t+g_t-L_t\Delta t-c_t-s_t\rvert\le10^{-6}$~kWh
       & 内：问题 1 为 $1.49\times10^{-13}$，问题 2、3 均为 $3.41\times10^{-13}$ & PASS \\
    A2 & 储电量递推残差 & 内 $\le10^{-6}$~kWh；交 $\le10^{-3}$~kWh
       & 内：$9.09\times10^{-13}$（问题 1）、$1.09\times10^{-11}$（问题 2、3）；交（问题 1、2、3、4-2、4-3）：$3.64\times10^{-12}$、$2.96\times10^{-4}$、$2.72\times10^{-4}$、$2.11\times10^{-4}$、$2.61\times10^{-4}$ & PASS \\
    A3 & 储电量越界数 & 等于 $0$（运行区间 $[1200,\,10800]$~kWh）
       & 交：五个结果文件均为 $0$；内：问题 1--3 均为 $0$ & PASS \\
    A4 & 充放电功率越界数 & 等于 $0$（单槽上限 $\overline P\Delta t=833.3333$~kWh）
       & 交：块级上限 $4\overline P\Delta t=3333.3333$~kWh，四个横表均为 $0$；内：问题 1--3 的 10~min 级均为 $0$ & PASS \\
    A5 & 充放互斥违反数 & 等于 $0$（$\eta_c\eta_d=0.81<1$ 时最优解不同时充放）
       & 内：问题 1--3 均为 $0$，并在 97 个储电量双界均不活跃的槽位复核 & PASS \\
    A6 & 周期与终端条件 & 问题 1 的 $\lvert S_0-S_{144}\rvert=0$；问题 2--4 每日 $S_{144}\ge S_0$
       & 问题 1 的日周期偏差 $0.0$~kWh；四个横表「末值低于初值的天数」均为 $0$ & PASS \\
    A7 & 非负性违反数 & 等于 $0$
       & 交：负购电量 $0$ 个（四个横表各 $334\times144$ 个时段值，加 result1 的 144 个）；内：问题 1--3 均为 $0$ & PASS \\
    A8 & 费用复算相对误差 & $\le10^{-4}$
       & 交：问题 1 为 $0.0$；内：问题 1--3 均为 $0.0$ & PASS \\
    A9 & 日期与结构 & 334 天；工作表名、列数与列序逐项一致
       & 四个横表均为 $334\times147$，日期 2025-02-01 至 2025-12-31；result1 为 144 行计划表加 6 行充放电表；工作表名 result1 为计划购电量与充放电量，result2、result4-2 为计划购电量、充放电量、紧急购电量，result3、result4-3 为计划购电量、调整购电量、充放电量、紧急购电量 & PASS \\
    A10 & 全天汇总列 & 购电量偏差 $\le10^{-2}$~kWh（144 槽取整上界 $7.2\times10^{-3}$）；费用相对偏差 $\le10^{-4}$
       & 购电量最大偏差 $1.10\times10^{-3}$、$1.10\times10^{-3}$、$1.00\times10^{-3}$、$1.20\times10^{-3}$~kWh（对应问题 2、3、4-2、4-3）；费用最大相对偏差 $3.40\times10^{-7}$、$3.40\times10^{-7}$、$4.32\times10^{-7}$、$4.05\times10^{-7}$ & PASS \\
    \bottomrule
  \end{tabularx}
\end{table}

\paragraph{反面证据：结算平衡不得漏掉紧急购电项}

附件 5 把「计划购电量」与「紧急购电量」分成\textbf{两张表}，因此计划量不等于实际供给量。
若 A1 的结算平衡漏掉紧急购电项 $g_t$（即把计划购电量当作全部供给），残差恰好等于 $g_t$：
在问题 2 的 334 天逐槽解上 $\max_t g_t=326.18$~kWh，在问题 3 上为 $204.87$~kWh，
分别比 $10^{-6}$~kWh 的容差大 $3.3\times10^{8}$ 倍与 $2.0\times10^{8}$ 倍；
补上 $+g_t$ 后，问题 2、3 的平衡残差降为 $3.41\times10^{-13}$~kWh。
同一次审计还给出：问题 2 的盈余处置量（弃光与超购部分的处置量）为 $3\,106\,390.84$~kWh，
问题 3 为 $2\,638\,963.35$~kWh，均为正残差，与「供给不低于需求」的口径一致。

\paragraph{审计结论}

\begin{center}
\fbox{\ 全部通过，可进入写作\ }\qquad
\fbox{\ 存在失败项，已定位并修正\ }（未勾选）
\end{center}

\noindent A1--A10 十项全部通过；审计同时复核了交付物的结构与全天汇总列，
未发现任何越界、违反非负性或费用复算不一致。

\section{源程序清单}\label{sec:app-code}

以下为本论文建模、求解、绘图、审计与数值汇总所用的\textbf{全部源程序}（格式规范第五条），
与 \texttt{code/} 下的实际文件逐项一致；附录用 \texttt{lstinputlisting} 原样引入文件内容。
\texttt{[plain]} 选项表示该脚本不参与求解，只做绘图或文本处理。

\codefile{config.py}
\codefile{lp_engine.py}
\codefile{q1.py}
\codefile{q2.py}
\codefile{q3.py}
\codefile{q4.py}
\codefile{results_io.py}
\codefile{exp_scan.py}
\codefile{exp_sobol.py}
\codefile{exp_na.py}
\codefile{exp_oos.py}
\codefile{exp_benchmark.py}
\codefile{audit.py}
\codefile{paper_numbers.py}
\codefile{run_all.py}
\codefile{requirements.txt}
\codefile[plain]{plot_figures.py}
\codefile[plain]{plot_figures_v2.py}
\codefile{copula.py}
\codefile{scen_copula.py}
\codefile{exp_copula.py}
\codefile{exp_chance.py}
\codefile{exp_dro.py}
\codefile{decomp.py}

\begin{guideblock}
附录放\textbf{可复制的纯文本源码}，\textbf{不要用截图}，也不要放空壳函数。
每张数据图都要能追到生成脚本：图 3--图 10 由 \texttt{plot\_figures.py} 一次生成（每张图一个函数），
图 1、图 2 为 \texttt{figures/} 下的 TikZ 源文件、由 XeLaTeX 直接排版。
新增脚本必须同步补入本清单，否则「论文未列全部源程序」属格式规范第五条风险。
\end{guideblock}

\subsection{v2 新增实验的中间结果与复算入口}

v2 在原有主模型之外新增四组对照实验与两组拟合，全部数字同样只由脚本产生；
下表给出每个中间结果文件的复算入口与在正文中的引用位置。

\begin{table}[htbp]
  \centering
  \caption{v2 新增中间结果文件的复算入口（全部位于 \texttt{intermediate/} 目录）}
  \label{tab:v2map}
  \small
  \begin{tabular}{lll}
    \toprule
    中间结果文件 & 生成脚本 & 正文引用位置 \\
    \midrule
    \texttt{exp\_copula.json} & \texttt{exp\_copula.py}
      & \bref{sec:vine}（C 藤选族）、\ref{sec:q4-gen}（生成器对照）、\ref{sec:q4-risk}（分时段风险画像） \\
    \texttt{exp\_chance.json} & \texttt{exp\_chance.py} & \bref{sec:q2-chance}（机会约束对照与报童互证） \\
    \texttt{exp\_dro.json} & \texttt{exp\_dro.py} & \ref{sec:q4-dro}（盒式分布鲁棒的保守度—费用曲线） \\
    \texttt{decomp.json} & \texttt{decomp.py} & \ref{sec:lshape}（L 形分解与直接求解对拍） \\
    \bottomrule
  \end{tabular}
\end{table}

四组实验均复用 \texttt{config.py} 的口径（144 槽、$DT=1/6$~h、结果窗口 334 天、1 月预热、
终端 $S_{144}\ge S_0$、不允许售电），并各自包含独立自检：\texttt{exp\_dro.json} 的
\texttt{audit} 字段给出计划可行性、功率界、储电量界、递推残差、充放互斥与最坏情形覆盖的一致性残差；
\texttt{decomp.json} 的 \texttt{equivalence\_K60} 给出 334 天逐日的两算法目标值相对间隙；
\texttt{exp\_chance.json} 的 \texttt{cross\_check} 把两阶段随机规划的复算值与 \texttt{q2\_main.json}
对拍（相对差 $0.00$）；\texttt{exp\_copula.json} 的 \texttt{part3\_generators} 给出四种生成器的
样本外决策指标与边际保真度。

\section{数据字典}

\begin{table}[htbp]
  \centering
  \caption{附件字段、单位与时标约定}
  \label{tab:datadict}
  \footnotesize
  \begin{tabularx}{\linewidth}{C{1.4cm} L{3.0cm} C{1.6cm} C{2.4cm} Y}
    \toprule
    附件 & 字段 & 单位 & 时间范围 & 时标约定 \\
    \midrule
    附件 1 & 电价 & 元/kWh & 一天（144 时段） & 标签 = 时段末端；$0{:}00{+}1$ 表示次日 $0{:}00$（当日 $24{:}00$） \\
    附件 1 & 小区负载 & kW & 同上 & 同上，\textbf{实为附件 2 全年逐时段均值}（见下） \\
    附件 1 & 光伏发电「预测」功率 & kW & 同上 & \textbf{实为附件 2 全年逐时段均值，不是预报}（见下） \\
    附件 2 & 小区负载 & kW & 2025-01-01 至 12-31 & 365 行 $\times$ 144 列，标签 = 时段末端 \\
    附件 2 & 光伏发电实际功率 & kW & 同上 & 同上 \\
    附件 3 & 光伏发电预报 & kW & 同上 & 每天 $0/6/12/18$ 时发布，预报 $k$ 小时对应 $(h_0+k-1,\ h_0+k]$ \\
    附件 4 & 电价 & 元/kWh & 同上 & 365 行 $\times$ 144 列，标签 = 时段末端 \\
    附件 5 & 结果文件模板 & --- & 同上 & 工作表名、列名与列序即本文结果文件的格式约定 \\
    \bottomrule
  \end{tabularx}
\end{table}

\paragraph{附件 1 的身份核查（数据字典的关键一条）}

附件 1 的\textbf{电价、小区负载、光伏发电「预测」功率三列，实为附件 2 与附件 4 的全年逐时段均值，
并不是日前预报}：以 365 天逐时段均值逐槽比对，电价的最大偏差为 $5.15\times10^{-5}$~元/kWh
（\texttt{intermediate/q4\_price.json} 的 \texttt{skeleton\_is\_mean} 字段），
负载的最大偏差为 $1\times10^{-4}$~kW（\texttt{paper\_numbers.json}），光伏的最大偏差为 0.0431~kW。
因此附件 1 \textbf{只能用于问题 1 与全文基线口径}：问题 1 的典型日与无储能基线 $C^{base}$ 用它，
问题 2--4 的日前计划一律改用附件 3 的当次预报（问题 2 用 $0{:}00$ 预报，问题 3 在
$0{:}00$、$6{:}00$、$12{:}00$、$18{:}00$ 四次更新），结算一律用附件 2 与附件 4 的实际值。
若把附件 1 当作日前预报，计划就脱离了该时刻真正可得的信息，属信息结构的口径错误。
另有一个窗口口径需同时说明：只用结果窗口 2025-02-01 至 12-31 的 334 天求均值时，
电价的最大偏差约 $3.1\times10^{-2}$~元/kWh，即\textbf{窗口内均值并不完全等于附件 1}，
故问题 2--4 也不能把附件 1 当作窗口内的均值曲线使用。

（符号说明表已移入正文第 2.3 节，见正文表~\bref{tab:notation}。）


\subsection{相依结构：C 藤 pair-Copula 的逐层选族与 AIC 对照}\label{sec:vine-app}

正文 \S3.7 给出三点读数与结论；本小节给出逐层选族结果（\ref{tab:vine}）与五族候选的完整 AIC
对照（\ref{tab:vine-cand}），供复核。

\begin{table}[htbp]
  \centering
  \caption{C 藤 pair-Copula 逐层选族结果（样本为 2025-02-01 至 12-31 的 334 个日水平观测；
           AIC 越小越好，经验距离为与经验 Copula 的 \(L_2\) 距离）}
  \label{tab:vine}
  \footnotesize
  \setlength{\tabcolsep}{4pt}
  \begin{tabularx}{\linewidth}{L{2.4cm} C{1.2cm} L{2.55cm} C{1.15cm} C{1.15cm} C{1.15cm} Y}
    \toprule
    变量对 & 最优族 & 参数 & \(\tau\) & \(\lambda_U\) & \(\lambda_L\) & AIC / 经验距离 \\
    \midrule
    电价—负载 & Frank & \(\theta=10.4686\) & \(+0.6976\) & 0 & 0 & \(-453.54\) / \(0.0135\) \\
    电价—光伏 & \(t\) & \(\rho=-0.2850,\ \nu=3.2324\) & \(-0.2383\) & \(0.0479\) & \(0.0479\)
      & \(-44.50\) / \(0.0210\) \\
    负载—光伏 \(|\) 电价 & Frank & \(\theta=6.3107\) & \(+0.5542\) & 0 & 0 & \(-253.44\) / \(0.0281\) \\
    \bottomrule
  \end{tabularx}
\end{table}

\textbf{三点读数。} ① \textbf{电价与负载强正相依}（\(\tau=+0.698\)、Pearson \(+0.919\)），
但两端\textbf{无尾相依}（\(\lambda_U=\lambda_L=0\)）：贵的时候确实多用电，但这种关系在尾部并不额外
放大。② \textbf{电价与光伏为弱负相依并带轻微对称尾相依}：最优族是自由度仅 3.23 的 \(t\) Copula，
\(\tau=-0.238\)、\(\lambda_U=\lambda_L=0.0479\)。这正是“光伏极低的日子电价极高”这类\textbf{尾部
共现}的量化结果——它存在，但强度只有约 5\%，既不能忽略也不应夸大；若改用 Gaussian 族（AIC 高 14.89，
且强制 \(\lambda=0\)）就会把它抹平。③ \textbf{控制电价后负载与光伏仍强正相依}（\(\tau=+0.554\)），
说明两者有共同驱动（季节与天气），并非仅是电价的代理变量。\evid{C-Q4-02}

五族候选的 AIC 对照（表~\ref{tab:vine-cand}）显示：三对变量的 AIC 最优族与经验 Copula 距离最优族
排序一致，且最优与次优的 AIC 差距（51.37、14.89、37.29）远大于拟合噪声，选族结论不依赖单一判据。

\begin{table}[htbp]
  \centering
  \caption{五族候选 Copula 的 AIC 对照（同一 334 个日水平观测；加粗为各族中的最小值，
           “\(|\)”后为条件变量）}
  \label{tab:vine-cand}
  \small
  \begin{tabular}{lccccc}
    \toprule
    变量对 & Gaussian & \(t\) & Clayton & Gumbel & Frank \\
    \midrule
    电价—负载 & \(-402.17\) & \(-398.97\) & \(-357.13\) & \(-305.67\) & $\mathbf{-453.54}$ \\
    电价—光伏 & \(-29.61\) & $\mathbf{-44.50}$ & $2.16$ & $2.22$ & \(-22.75\) \\
    负载—光伏 \(|\) 电价 & \(-216.15\) & \(-212.55\) & \(-153.59\) & \(-173.89\) & $\mathbf{-253.44}$ \\
    \bottomrule
  \end{tabular}
\end{table}


五族 AIC 对照见表~\ref{tab:vine-cand}。


\subsection{相关性三尺度与预报误差分档明细}\label{sec:datadetail-app}

正文 \S3.4 与 \S3.5 给出结论；本小节给出完整表格，供复核。

\begin{table}[htbp]
  \centering
  \caption{电价与负载、净负荷、光伏的相关系数（三个尺度；池化与去形状口径的样本量为
           \(365\times144=52\,560\) 个时段观测，日级口径为 365 个日观测）}
  \label{tab:corr-app}
  \small
  \begin{tabular}{lccc}
    \toprule
    变量对 & 池化（全部时段） & 日级（逐日先算再合并） & 剔除日内形状后 \\
    \midrule
    电价 --- 负载 & \(+0.5077\) & \(+0.8977\) & \(+0.4554\) \\
    电价 --- 净负荷 \(L-PV\) & \(+0.2602\) & \(+0.9824\) & \(+0.1583\) \\
    电价 --- 光伏 & \(+0.0284\) & \(-0.3431\) & \(+0.0484\) \\
    \bottomrule
  \end{tabular}
\end{table}

\begin{table}[htbp]
  \centering
  \caption{光伏预报误差统计（样本为 2025 年每天 \(0{:}00\) 发布的 24~h 预报与附件 2 对应实际功率的
            逐槽配对，\(n=52\,560\)；MAE、RMSE、bias 单位均为 kW）}
  \label{tab:fcsterr-app}
  \footnotesize
  \setlength{\tabcolsep}{4pt}
  \begin{tabularx}{\linewidth}{L{3.1cm} C{1.5cm} C{1.5cm} C{1.5cm} C{1.5cm} Y}
    \toprule
    口径（目标整点范围） & 样本量 \(n\) & MAE & RMSE & bias & 结论 \\
    \midrule
    全时段（\(0\)--\(23\) 时，144 槽/天） & \(52\,560\) & \(361.58\) & \(636.37\) & \(-26.96\)
      & 夜间近零样本摊薄指标，只作对照 \\
    白天（\(\ge5\) 时，114 槽/天） & \(41\,610\) & \(456.72\) & \(715.22\) & \(-34.06\)
      & \textbf{建模依据}：MAE 为全时段的 1.26 倍 \\
    \midrule
    \multicolumn{6}{l}{\textbf{按目标整点聚合的偏差结构}（组内 bias 单位 kW）} \\
    \midrule
    夜间（光伏无有效出力） & \(0\)--\(4\) 时 & \multicolumn{2}{c}{\(0.00\) 至 \(+0.09\)}
      & \(4\) 时 \(+0.09\) & 无信息、无惩罚 \\
    上午（系统性\textbf{高估}） & \(5\)--\(10\) 时 & \multicolumn{2}{c}{\(+214.21\) 至 \(+733.00\)}
      & \(7\) 时 \(+733.00\) & 计划偏少 $\rightarrow$ 触发 \(5\) 倍紧急购电 \\
    正午过渡 & \(11\) 时 & \multicolumn{2}{c}{\(+7.21\)} & --- & 接近无偏 \\
    下午（系统性\textbf{低估}） & \(12\)--\(18\) 时 & \multicolumn{2}{c}{\(-801.87\) 至 \(-106.46\)}
      & \(15\) 时 \(-801.87\) & 计划偏多 $\rightarrow$ 白付电费或付违约电价 \\
    傍晚之后（预报空洞） & \(19\)--\(23\) 时 & \multicolumn{2}{c}{\(-0.15\) 至 \(0.00\)}
      & \(19\) 时 \(-0.15\) & 出力近零，不触发惩罚 \\
    \bottomrule
  \end{tabularx}
\end{table}

\subsection{联合情景保真度检验明细}\label{sec:scenfid-app}

正文 \S7.2 给出结论；本小节给出完整检验表。

\begin{table}[htbp]
  \centering
  \caption{联合情景的保真度检验（对象：64 个代表历史日与 64 个分块 Bootstrap 情景；参照：2025-02-01 至
           12-31 的 334 天实际数据；相对偏差单位 \%，相关系数为无量纲量）}
  \label{tab:scenfid}
  \footnotesize
  \setlength{\tabcolsep}{4pt}
  \begin{tabularx}{\linewidth}{L{4.35cm} C{2.35cm} C{2.35cm} Y}
    \toprule
    检验指标 & 代表历史日（主口径） & 分块 Bootstrap（对照） & 判读 \\
    \midrule
    电价均值相对偏差 & 0.08\% & 3.41\% & 主口径的均值几乎无偏 \\
    最大分位相对差（7 个分位点） & 6.99\%（1\% 分位） & 20.4\%（1\% 分位）
      & 两者都在低价尾部误差最大，对照口径的尾部偏差约为主口径的 3 倍 \\
    情景 \(\rho(p,\text{净负荷})\) & 0.2597（实际 0.2620） & 无法按日配对
      & 主口径 \textbf{保留}了联合结构 \\
    日净需求残差均值 & 保持 \(-715.85\)~kWh（全样本 \(-712.23\)） & 不适用（块级拼装）
      & 净需求侧边缘分布同样保真 \\
    \bottomrule
  \end{tabularx}
\end{table}
\section{对照实验明细}\label{sec:ctl-app}

正文 \S7.3 的三组对照实验只保留结论；本节给出其完整表格、逐条数据与判定阈值，便于复核。
三个小节沿用正文的小节标题与编号，便于互查。

\subsection{情景生成方式的样本外对照：结论对生成器的敏感性}
\label{sec:q4-gen}

\bref{sec:q4-price} 的联合情景用“二维分层取代表历史日”构造。这一做法是否只是“拟合了样本内的 334 天”？
换一种生成方式，结论会不会变？本节用\textbf{样本外对照实验}回答：情景集只用 \textbf{2025-02-01 至 06-30}
（150 天）拟合，在 \textbf{07-01 至 12-31}（184 天）上按真实轨迹结算，四种生成器共用同一决策模型、
同一价格信息结构（假设 A10）与同一终端条件：\textbf{（1）历史代表日}（主方案，二维分层 \(+\) 层内
medoid）；\textbf{（2）分块 Bootstrap}（6~h 块按时刻对齐重采样）；\textbf{（3）KDE \(+\) Copula}
（日水平核密度边缘 \(+\) 最优二元 Copula，抽样 600 次后削减到与（1）相同的情景数）；
\textbf{（4）按光伏规模缩放}（电价情景同（1），残差情景改为“训练窗相对残差形状 \(\times\) 当日光伏
预报电量”）。

\begin{table}[htbp]
  \centering
  \caption{四种情景生成器的样本外对照（训练窗 2025-02-01 至 06-30 共 150 天，测试窗 07-01 至 12-31 共 184 天；
           同一决策模型、同一实际轨迹；费用单位元，电量为 184 天累计 kWh，“边际偏差”为日总残差
           1\%/99\% 分位相对训练窗实际分布的最大绝对偏差）}
  \label{tab:gen-cmp}
  \small
  \begin{tabular}{lrrrrr}
    \toprule
    生成器 & 总费用 & 相对（1） & 紧急购电量 & 逐日 CVaR$_{0.95}$ & 边际偏差 \\
    \midrule
    （1）历史代表日 & $9\,046\,598.58$ & — & $76\,033.90$ & $79\,378.60$ & $345.7$ \\
    （2）分块 Bootstrap & $9\,038\,928.08$ & $-0.085\%$ & $72\,324.35$ & $79\,381.24$ & $2\,102.3$ \\
    （3）KDE $+$ Copula & $9\,072\,284.06$ & $+0.284\%$ & $107\,458.94$ & $80\,061.08$ & $1\,519.1$ \\
    （4）按光伏规模缩放 & $9\,068\,058.80$ & $+0.237\%$ & $92\,727.06$ & $81\,776.72$ & 不适用 \\
    \bottomrule
  \end{tabular}
\end{table}

\textbf{三条结论。} ① \textbf{费用结论对生成器不敏感}：四种生成器的样本外总费用极差
33\,355.98 元，仅为四者均值（9\,056\,467.38 元）的 0.37\%，相对无储能基线的节省率落在
11.10\%--11.43\%。这比“样本外劣化”检验更强：不仅主方案不过拟合，连换一种生成范式也不改变结论。
② \textbf{但分项指标并不等价}：紧急购电量在四种生成器下为 72\,324--107\,459~kWh（相差 48.5\%）、
逐日 CVaR\(_{0.95}\) 为 79\,379--81\,777~元（相差 3.0\%）——\textbf{费用稳健，而“靠多少紧急购电兜底”
可被生成方式改变}，引用该分项必须同时给出情景口径。③ \textbf{联合尾部的保真度上 Copula 更好}：
“高价 \(+\) 光伏高估”的联合尾共现率真值 0.0533，Copula 生成器 0.0500、历史代表日 0.0667，分块
Bootstrap 只有 0.0167；这与 \bref{sec:vine} 一致——尾部共现在数据中\textbf{本来就弱}（\(\lambda\approx0.048\)），
其保真度差异不足以传导到决策层。据此本文\textbf{主模型保留历史代表日生成器}，\textbf{Copula 用于相依
结构诊断}（\bref{sec:vine}）与\textbf{敏感性对照}（本小节），两者都不改变任何主结果。\evid{C-Q4-03}



\subsection{分时段 CVaR 风险画像}
\label{sec:q4-risk}

“全年总费用”只有一个数，掩盖了风险的时段结构（分时段风险度量在电力市场风险管理中已有应用~\cite{refwangtaiqiang}）。本节把费用\textbf{按槽位分解}：
第 $d$ 天第 $t$ 槽的费用为 $c_{d,t}=p_{d,t}E_{d,t}+5\,p_{d,t}g_{d,t}$（计划购电 $+$ 紧急购电），
对固定的 $t$ 取跨 334 天的 $\mathrm{CVaR}_{0.95}$，再按小时求和，得到\textbf{逐小时风险画像}
（图~\ref{fig:risk-profile}；问题 2 与问题 4-2 分别计算）。

\begin{figure}[htbp]
  \centering
  \includefig[width=0.68\linewidth]{fig13_risk_profile.pdf}
  \caption{分时段 CVaR 风险画像（柱：各槽位费用的 \(\mathrm{CVaR}_{0.95}\) 按小时求和；实线：逐小时均值；
            虚线：槽位缺电频率（右轴）；左：问题 2，右：问题 4-2；范围：334 天 \(\times\) 144 槽）}
  \label{fig:risk-profile}
\end{figure}

\textbf{风险的时段分化是结构性的。} ① \textbf{费用尾部集中在傍晚}：固定价下逐小时
\(\mathrm{CVaR}_{0.95}\) 的高点出现在 \(15\)--\(18\) 时（6\,075、5\,566、6\,604、6\,358 元/日），
波动价下同段升到 7\,492、6\,778、7\,550、8\,098 元/日。② \textbf{缺电尾部集中在上午}：槽位缺电频率
在 \(5\)--\(11\) 时高达 14.5\%--23.8\%（固定价）与 14.8\%--34.6\%（波动价），峰值分别在 \(11\)
时（23.8\%）与 \(5\) 时（34.6\%），而 \(0\)--\(3\) 时与 \(19\)--\(23\) 时几乎为 0——这与
\S3.5 的偏差剖面完全对应（上午光伏被系统性高估使 \(0{:}00\) 计划偏少）。③ \textbf{风险集中度}：
\(\mathrm{CVaR}\) 最高的 12 个槽位只占全部槽位之和的 16.73\%（问题 2）与 16.10\%（问题 4-2），
说明风险不是集中在个别槽位而是沿上午与傍晚两个时段带分布；逐槽 \(\mathrm{CVaR}\) 之和是逐槽均值之和的
2.01 倍（固定价）与 2.34 倍（波动价），即\textbf{电价波动把费用的尾部放大得比均值多}。\evid{C-Q4-04}

\textbf{决策含义。} 安全边际应当\textbf{按目标整点}配置（上午与傍晚需要，\(0\)--\(4\) 时不需要），
而不是全天统一加比例；滚动调整的价值同样集中在这两个时段带。


\subsection{分布鲁棒对照：保守度—费用曲线}
\label{sec:q4-dro}

\ref{sec:q4-gen} 的四种生成器都属\textbf{随机规划}：先估计一个分布，再在期望意义下优化。当样本有限或分布可能
误设时，另一种思路是\textbf{分布鲁棒优化}（DRO）——在模糊集上取最坏期望，其最简单的成员就是
\textbf{盒式不确定集}（矩约束模糊集下的多阶段建模见文献~\cite{refyuhaodong}）。本节把它作为
\textbf{保守度对照}：把“保守度”参数化后定量回答一个问题——\textbf{用多少实际费用，能买到多少最坏
情形下的改善？}

\textbf{模型。} 取逐槽盒式不确定集
\(p_t\in[\bar p_t-k\sigma^p_t,\ \bar p_t+k\sigma^p_t]\)、
\(\widetilde{ND}_t\in[ND^{\text{fcst}}_t-k\sigma^{nd}_t,\ ND^{\text{fcst}}_t+k\sigma^{nd}_t]\)，其中 \(\sigma^p_t\) 与 \(\sigma^{nd}_t\) 分别为附件 4 电价的逐槽样本
标准差与 \(0{:}00\) 预报残差的逐槽样本标准差（365 天，\(\mathrm{ddof}=1\)）。目标在盒的顶点取得，最坏情形目标与约束取
\(\min\ \sum_t p^{\text{hi}}_tE_t+5\sum_t p^{\text{hi}}_t z_t\)、\(z_t\ge ND^{\text{hi}}_t+c_t-d_t-E_t\)、\(z_t\ge0\)，外加计划可行性与储电量约束（同 \bref{sec:q4-price}）；\(k=0\) 时退化为\textbf{按预报的确定性购电}（零安全边际）。

\textbf{一个结构性推论（先于数值给出）。} 上述最坏情形目标中 $E_t$ 与 $z_t$ 在同槽位上的目标系数之比
恒为 $1:5$，而把 $E_t$ 提高 $\delta$ 可使 $z_t$ 下降同量、目标值下降 $4p^{\text{hi}}_t\delta>0$，
故只要 $z_t>0$ 就仍有改进空间：\textbf{最优解必取 $z_t\equiv0$}，即
$E_t=\max\{0,\ ND^{\text{hi}}_t+c_t-d_t\}$。也就是说，\textbf{盒式 DRO 在本问题设定下退化为
“按最坏参数做确定性优化”}——它逐槽覆盖 $k$ 倍标准差的不利偏差，而不依赖任何追索机制。
数值上 $z$ 的一致性残差为 $0.000$（表~\ref{tab:audit}），与推论一致。

\begin{table}[htbp]
  \centering
  \caption{盒式分布鲁棒对照：保守度 $k$ 与实际回测表现（334 天，2025-02-01 至 12-31；决策用盒式鲁棒模型，
           结算用附件 4 实际电价与附件 2 实际负载/光伏；$W_{k=3}$ 为各方案在\textbf{同一}
           $k=3$ 不确定集下的最坏情形目标值；费用单位元，电量为 kWh）}
  \label{tab:dro}
  \small
  \begin{tabular}{crrrrrr}
    \toprule
    $k$ & 回测总费用 & 节省率 & 紧急购电量 & 逐日 CVaR$_{0.95}$ & 最坏单日 & $W_{k=3}$（万元） \\
    \midrule
    $0.0$ & $18\,258\,473.08$ & $-7.23\%$ & $1\,156\,618.02$ & $91\,630.53$ & $101\,653.74$ & $5\,956.86$ \\
    $0.5$ & $17\,404\,407.87$ & $-2.22\%$ & $793\,522.04$ & $86\,473.81$ & $94\,613.68$ & $3\,557.47$ \\
    $1.0$ & $16\,936\,297.51$ & $+0.53\%$ & $521\,080.78$ & $84\,250.61$ & $89\,459.18$ & $3\,058.17$ \\
    $\mathbf{1.5}$ & $\mathbf{16\,819\,102.59}$ & $\mathbf{+1.22\%}$ & $328\,147.59$ & $84\,471.07$ & $87\,672.62$ & $2\,544.50$ \\
    $3.0$ & $18\,015\,766.37$ & $-5.81\%$ & $62\,776.53$ & $91\,718.61$ & $95\,526.83$ & $2\,512.55$ \\
    \midrule
    \textbf{随机规划}（\bref{sec:q4-price}、\bref{sec:q4-frontier}） & $\mathbf{15\,094\,255.84}$ & $\mathbf{+11.35\%}$ & $182\,759.86$ & $\mathbf{77\,697.98}$ & $\mathbf{81\,643.86}$ & — \\
    \bottomrule
  \end{tabular}
\end{table}

\textbf{四条结论。} ① \textbf{回测费用对 $k$ 呈 U 形}：$k$ 由 $0$ 增大到 $1.5$ 时，多买电的代价
小于少买电触发的 $5$ 倍紧急购电，费用由 $18\,258\,473.08$ 降到 $16\,819\,102.59$~元（省
$1\,439\,370.50$ 元）；此后计划购电量继续上升而紧急购电已接近清零，费用反弹，$k=3.0$ 时回到
$18\,015\,766.37$~元。实证最优点为 $k^{*}=1.5$。② \textbf{盒式鲁棒在最坏情形上确实有效}：在
\textbf{同一} $k=3$ 的不确定集下评估，$k=0$ 的方案最坏情形要花 $5\,956.86$ 万元，$k=3$ 的方案只要
$2\,512.55$ 万元（降 $57.8\%$）。③ \textbf{但它的“保守”买贵了}：从 $k^{*}$ 加到 $k=3.0$，实际回测
费用多付 $1\,196\,663.78$~元（$+7.12\%$），即最坏情形费用每降低 $1\%$、回测费用要多付 $3.98$ 万元
（约基线的 $0.234\%$）。④ \textbf{盒式鲁棒被随机规划全面占优}：与 \bref{sec:q4-frontier} 的随机规划相比，最优方案
（$k^{*}=1.5$）在\textbf{全部五项指标}上更差——总费用高 $1\,724\,846.75$~元（$+11.43\%$）、逐日
$\mathrm{CVaR}_{0.95}$ 高 $6\,773.09$~元/日（$+8.72\%$）、最坏单日高 $6\,028.76$~元（$+7.38\%$）、
紧急购电量多 $145\,387.73$~kWh（$+79.6\%$）、槽位缺电频率反而更高（$10.24\%$ 对 $9.04\%$）。
\evid{C-Q4-05}

\textbf{为什么会这样，以及正确的使用方式。} 盒式不确定集把 $144$ 个槽位\textbf{同时}取到极端
（电价与净需求全线 $+k\sigma$），而实测 334 天里槽间残差并不同步（\S7.1 的日内自相关与 $19$--$23$ 时
近乎零的偏差都说明这一点），因此这个“最坏情形”几乎不可能发生，为它买保险的钱花得不值。此外 $k$ 的
经验最优值恰好使\textbf{电价加权的缺货频率}降到 $1/5$ 附近，与 \bref{sec:q2-newsvendor} 报童模型的允许缺货概率
$\beta^{*}=1/5=0.20$ 一致——两条完全不同的建模路线给出同一个服务水平。故盒式鲁棒的价值是
\textbf{压力上界}与\textbf{轻度鲁棒中间方案}（$k=0.5$--$1.0$），而非“更好的策略”；真正改进的方向是把
\textbf{全盒}换成\textbf{带不确定性预算 $\Gamma$ 的不确定集}或\textbf{矩约束模糊集}
\cite{refyuhaodong}，使“每个时段都可能极端、但不允许全部时段同时极端”这一先验进入模型（\S10.4）。



\section{口径定义汇总}\label{sec:caliber-app}

正文各章在本题参数范围内用到时以~\calref{C-0x} 引用本节的对应条目。口径编号在建模前固定、未随结果调整；
每个口径含定义、公式、数据来源、时段、单位、不含项、对照与复算入口八项。

g

g

g

g

g

\section{长推导}

\subsection{对偶推导与阈值律}

记式~\bref{eq:q1-balance} 的乘子为 $\alpha_t$（电量平衡的影子价格，
在正常购电时段等于电价，见式~\eqref{eq:app-alpha-eq-p}），式~\bref{eq:q1-soc} 的乘子为
$\nu_t$（储电量的边际价值，即水价值），储电量上下界乘子为 $\lambda^{\pm}_t\ge 0$，
充放电功率界乘子为 $\varrho^{c}_t,\varrho^{d}_t\ge 0$。
把电量平衡写成 $(L_t-PV_t)\Delta t+c_t+s_t-E_t-d_t=0$，拉格朗日函数为

\begin{equation}
  \begin{aligned}
    \mathcal L=\ &\sum_{t}p_tE_t
    +\sum_{t}\alpha_t\Bigl[(L_t-PV_t)\Delta t+c_t+s_t-E_t-d_t\Bigr]\\
    &+\sum_{t}\nu_t\Bigl[S_t-S_{t-1}-\eta_cc_t+\frac{d_t}{\eta_d}\Bigr]
    +\sum_{t}\lambda^{+}_t\bigl[S_t-\overline S\bigr]
    +\sum_{t}\lambda^{-}_t\bigl[\underline S-S_t\bigr]\\
    &+\sum_{t}\varrho^{c}_t\bigl[c_t-\overline P\Delta t\bigr]
    +\sum_{t}\varrho^{d}_t\bigl[d_t-\overline P\Delta t\bigr]
  \end{aligned}
  \label{eq:app-lag}
\end{equation}

对各决策变量求平稳性条件（$E_t,s_t,c_t,d_t$ 对 $t=1,\dots,144$；$S_t$ 对 $t=1,\dots,143$，
因为日周期约束把 $S_{144}=S_0=6000$~kWh 固定）：

\begin{align}
  \frac{\partial\mathcal L}{\partial E_t}&=p_t-\alpha_t\ \ge 0,
  &E_t\bigl(p_t-\alpha_t\bigr)&=0 \label{eq:app-stat-E}\\
  \frac{\partial\mathcal L}{\partial s_t}&=\alpha_t\ \ge 0,
  &s_t\,\alpha_t&=0 \label{eq:app-stat-s}\\
  \frac{\partial\mathcal L}{\partial c_t}&=\alpha_t-\eta_c\nu_t+\varrho^{c}_t\ \ge 0,
  &c_t\bigl(\alpha_t-\eta_c\nu_t+\varrho^{c}_t\bigr)&=0 \label{eq:app-stat-c}\\
  \frac{\partial\mathcal L}{\partial d_t}&=-\alpha_t+\frac{\nu_t}{\eta_d}+\varrho^{d}_t\ \ge 0,
  &d_t\Bigl(-\alpha_t+\frac{\nu_t}{\eta_d}+\varrho^{d}_t\Bigr)&=0 \label{eq:app-stat-d}\\
  \frac{\partial\mathcal L}{\partial S_t}&=\nu_t-\nu_{t+1}+\lambda^{+}_t-\lambda^{-}_t=0
  \label{eq:app-stat-S}
\end{align}

由式~\eqref{eq:app-stat-E}，购电时段（$E_t>0$）恒有

\begin{equation}
  \alpha_t=p_t
  \label{eq:app-alpha-eq-p}
\end{equation}

即电量平衡的影子价格等于该时段的电价；数值上逐槽复核的最大偏差为 $0.000$~元/kWh。
若充放电功率界均不活跃（$\varrho^{c}_t=\varrho^{d}_t=0$），式~\eqref{eq:app-stat-c} 与
式~\eqref{eq:app-stat-d} 化为

\begin{equation}
  \alpha_t-\eta_c\nu_t\ \ge 0,\qquad
  \frac{\nu_t}{\eta_d}-\alpha_t\ \ge 0
  \label{eq:app-threshold-raw}
\end{equation}

且 $c_t>0$ 时前式取等、$d_t>0$ 时后式取等，即得正文命题~\bref{prop:threshold} 的阈值律：
\textbf{$\alpha_t\le\eta_c\nu_t$ 时充电、$\alpha_t\ge\nu_t/\eta_d$ 时放电}，
死区为 $\eta_c\nu_t<\alpha_t<\nu_t/\eta_d$，死区比为 $1/(\eta_c\eta_d)=1.2346$
（单边效率各 0.9、往返 0.81），与式~\bref{eq:q1-deadband} 一致。
由式~\eqref{eq:app-stat-S} 还可得水价值的跨期关系：在储电量上下界均不活跃的连续时段上
$\nu_t=\nu_{t+1}$（水价值恒定）；储电量触及上界时 $\lambda^{+}_t>0$、$\nu_t<\nu_{t+1}$；
触及下界时 $\lambda^{-}_t>0$、$\nu_t>\nu_{t+1}$。
这一关系是图~\bref{fig:dualprice} 中水价值曲线出现折点的原因。

数值复核（\texttt{intermediate/q1\_dual2.json}）：储电量双界均不活跃的槽位 97 个；
$\nu_t\in[0.4711,\ 1.1305]$~元/kWh，$\alpha_t\in[0.3713,\ 1.2561]$~元/kWh；
全部 144 槽的充放电阈值律违反次数为 $0$，互补松弛最大残差为 $0.000$~元/kWh，
式~\eqref{eq:app-alpha-eq-p} 的最大偏差为 $0.000$~元/kWh。

\subsection{报童临界分位数的推导}

设时段 $t$ 的需求 $D_t$ 为随机变量、分布函数为 $F_t$，计划购电量为 $E_t$，
过剩成本为 $c_o$（多买 1~kWh 的损失），缺货成本为 $c_u$（不足 1~kWh 的额外支出）。
期望总成本为

\begin{equation}
  \Pi(E)=c_o\,\mathbb E\bigl[(E-D)^{+}\bigr]+c_u\,\mathbb E\bigl[(D-E)^{+}\bigr]
  \label{eq:app-nv-cost}
\end{equation}

对 $E$ 求导并令其为零：

\begin{equation}
  \frac{\mathrm d\Pi}{\mathrm dE}
  =c_oF(E)-c_u\bigl[1-F(E)\bigr]=0
  \quad\Longrightarrow\quad
  F(E^{*})=\frac{c_u}{c_u+c_o}
  \label{eq:app-nv-foc}
\end{equation}

本题中不足 1~kWh 需按 $5p_t$ 紧急购电，而正常电价 $p_t$ 已随计划购电量支付，
故缺货的额外支出为 $c_u=5p_t-p_t=4p_t$；超购的电量若无处可用（超购即弃置），
则 $c_o=p_t$，代入式~\eqref{eq:app-nv-foc} 得

\begin{equation}
  F(E^{*})=\frac{4p_t}{4p_t+p_t}=\frac45=0.8
  \qquad\Longrightarrow\qquad
  \text{最优缺货概率}=1-F(E^{*})=\frac15=0.2
  \label{eq:app-nv-result}
\end{equation}

即「留安全边际」的临界分位数为 $4/5$、最优缺货概率为 $1/5$。
若超购电量能转入储能、留待后续更高价时段使用（单位回收价值 $r_t$），
则 $c_o=p_t-r_t$，临界分位数相应上移：取全年均价 $p=0.7662$~元/kWh 与全年最低价
$p^{low}=0.3713$~元/kWh，得 $c_o\approx0.395$~元/kWh，$F(E^{*})\approx0.886$，
最优缺货概率降至约 $1/10$。此即正文式~\bref{eq:q2-newsvendor} 与表~\bref{tab:newsvendor} 的来源。

实测对照（\texttt{paper\_numbers.json} 的「报童经验」）：主口径下逐槽缺货频率为 $8.55\%$，
缺口电量 $172\,015.82$~kWh，电量满足率 $99.04\%$，安全边际合计 $576\,299.67$~kWh
（$11.98$~kWh/槽），落在 $1/10$ 与 $1/5$ 两界\textbf{之下侧}——
因为储能同时承担了对冲功能，报童结论只作方向性近似与上下界使用，不作精确刻画。

\section{随机种子清单}

本工程\textbf{只在四处使用随机数}，四问的线性规划求解、主口径的场景削减、
审计与数值汇总均为确定性过程（同一输入必得同一输出）。逐项清单如表~\ref{tab:seeds}。

\begin{table}[htbp]
  \centering
  \caption{随机种子清单（代码中的全部随机源；其余过程确定）}
  \label{tab:seeds}
  \footnotesize
  \setlength{\tabcolsep}{4pt}
  \begin{tabularx}{\linewidth}{L{3.0cm} L{3.2cm} C{2.5cm} Y}
    \toprule
    用途 & 随机机制 & 种子 & 说明 \\
    \midrule
    问题 4 联合情景合成（条件对照） & 分块 Bootstrap，按时刻对齐 6~h 块重采样
      （\texttt{q4.build\_joint\_scenarios(mode=“boot”)}） & \texttt{numpy.random.default\_rng(seed=7)}
      & 电价块与光伏预报残差块成对抽取，保住日内形状与相关性；一次生成后固定使用 \\
    场景削减的对照路径 & $k$-means 初始化（\texttt{lp\_engine.reduce\_scenarios}）
      & \texttt{random\_state=0}（\texttt{n\_init=4}）
      & 仅在 \texttt{--scen kmeans} 对照下启用；\textbf{主口径不抽样}：按日总残差等频分层后在层内取 medoid，确定性 \\
    非预期性置换检验 & 对某次预报做日间随机置换（\texttt{exp\_na.py}）
      & \texttt{numpy.random.default\_rng(seed=0)}；对照用 \texttt{seed=1}
      & 检验「置换 $6{:}00$ 预报后槽 0--35 完全不变、槽 36--71 改变」，用于排除信息泄漏 \\
    全局灵敏度的采样 & Sobol 低差异序列置乱
      （\texttt{scipy.stats.qmc.Sobol(scramble=True)}，\texttt{exp\_sobol.py}）
      & 费用目标 \texttt{seed=42}；节省率目标 \texttt{seed=43}
      & 置乱只改变采样点的取法，不改变参数区间与模型本身 \\
    \bottomrule
  \end{tabularx}
\end{table}

\noindent\textbf{确定性的部分（无随机数）}：问题 1--4 的全部线性规划由 SciPy
\texttt{linprog(method=“highs”)} 精确求解；问题 2、3 的场景削减用「等频分层加层内 medoid」，
问题 4 的历史代表日情景用「日净需求残差与日均电价」二维分层取 medoid、按层样本占比赋权，
两者都是确定性算法；\texttt{audit.py}、\texttt{paper\_numbers.py} 与 \texttt{exp\_scan.py}
不含随机数。因此论文的每一个数字在同一环境下重跑即可复现。

\paragraph{诚实说明}：本工程\textbf{未开展「多种子重复生成情景」的稳健性实验}，
故不报告费用的均值、标准差与 $95\%$ 置信区间，也不列出 20 个种子。
本工程对这一缺口的处理是：把随机性限制在只影响未来情景、不影响已实现时段决策的合成与采样环节
（分块 Bootstrap 只作为问题 4 联合情景的条件对照，主口径用确定性的分层 medoid），
并以同源对照替代多种子重复——同一固定情景集下的信息结构对照（R1、R2、R3）、
场景数对照（$K=60$ 与 $K=120$）、终端条件对照与安全边际消融，四者的结论方向一致。
若需要多种子区间，只需改动 \texttt{build\_joint\_scenarios} 的 \texttt{seed} 参数并重跑
\texttt{run\_all.py}，其余环节无需修改。

\section{验证与灵敏度明细}

本节汇集正文第九章与第三章各表的完整明细，便于复核；正文已给出各表的结论与判据。

\subsection{非预期性（信息泄漏）置换检验的逐区明细}\label{sec:na-app}

把某一次发布的预报在 12 天之间随机置换后重解，观察四个控制区（槽 $0$--$36$、$36$--$72$、
$72$--$108$、$108$--$144$ 分别由 $0{:}00$、$6{:}00$、$12{:}00$、$18{:}00$ 预报控制）的最大变化量，
结果见表~\ref{tab:na-app}。

\begin{table}[htbp]
  \centering
  \caption{非预期性置换检验：置换某一个发布时刻的预报后，四个控制区的最大变化量
           （对象：单日两阶段模型；范围：四个控制区各 36 个 10~min 槽；单位 kWh；
           变化恰为 $0.000$ 表示该区决策不受该次预报影响）}
  \label{tab:na-app}
  \footnotesize
  \setlength{\tabcolsep}{5pt}
  \begin{tabular}{lcccc}
    \toprule
    被置换的预报 & 控制区 1（$0{:}00$--$6{:}00$） & 控制区 2（$6{:}00$--$12{:}00$）
                 & 控制区 3（$12{:}00$--$18{:}00$） & 控制区 4（$18{:}00$--$24{:}00$） \\
    \midrule
    \(0{:}00\) & 833.33 & 822.06 & 489.52 & 0 \\
    \(6{:}00\) & 0.000  & 858.80 & 833.33 & 0 \\
    \(12{:}00\)& 0.000  & 0.000  & 805.35 & 0 \\
    \(18{:}00\)& 0.000  & 0.000  & 0.000  & 0 \\
    \bottomrule
  \end{tabular}
  \\[1ex]
  {\footnotesize 判读：任何一次预报被置换，都不改变早于其发布时刻的决策（变化恰为 $0.000$），
   故不存在用未来信息决策的泄漏；而“置换 $0{:}00$ 预报会改变全天决策”说明检验本身有效力。
   $18{:}00$ 预报的零影响源于其覆盖的傍晚时段光伏实际与预报都接近 0、信息量极低。}
\end{table}


\subsection{电价过程的四项统计检验：完整统计量与判定}\label{sec:price-test-app}

正文 \S7.1 给出三条建模含义；本小节给出四项检验（T1--T4）的完整统计量、结果与逐项判定。

\begin{table}[htbp]
  \centering
  \caption{电价过程的四项统计检验（对象：附件 4 的逐时段电价与残差过程；范围：2025-01-01 至 12-31、
            逐 10~min 共 \(365\times144=52\,560\) 个观测；单位为量纲一的统计量）}
  \label{tab:price-test}
  \footnotesize
  \setlength{\tabcolsep}{3pt}
  \begin{tabularx}{\linewidth}{C{0.62cm} L{1.85cm} L{4.35cm} L{3.35cm} Y}
    \toprule
    编号 & 检验对象 & 检验方法与统计量 & 结果 & 对建模的含义 \\
    \midrule
    T1 & 确定性日内骨架是否存在
      & 逐时段均值与逐时段标准差的对比（变异系数）
      & 附件 1 剖面与附件 4 逐时段均值的最大偏差 \(5.15\times10^{-5}\)~元/kWh（全年 365 天口径）；
        残差逐时段标准差均值 0.1238~元/kWh，变异系数 16.94\%
      & 骨架稳定存在但残差不可忽略，采用“骨架 \(+\) 残差”两层结构 \\
    T2 & 残差的正态性与日间差异
      & Jarque--Bera 检验、偏度与峰度；同一槽位跨日自相关
      & \(\mathrm{JB}=656.70\)（\(p=2.51\times10^{-143}\)）拒绝正态；偏度 \(-0.169\)、峰度 \(-0.431\)；
        日间自相关 \(\mathrm{ACF}(1\text{--}3)=0.513/0.148/0.144\)
      & 残差非正态、日间相关性弱：不做正态参数化，改用非参数分层重采样；骨架不按月分段 \\
    T3 & 日内自相关与波动结构
      & 残差自相关 \(\mathrm{ACF}(1\text{--}6)\)；峰段与谷段残差标准差之比
      & \(\mathrm{ACF}(1\text{--}6)\) 为 0.8589、0.666、0.614、0.573、0.494、0.432；峰段（19:00--20:00）
        残差标准差 0.1634~元/kWh 为谷段（03:00--04:00）0.0580~元/kWh 的 \textbf{2.81 倍}
      & 不能按独立同分布抽样：情景须按时刻对齐分块，同时保留日内相关与按时段异方差 \\
    T4 & 与负载、光伏的相关性
      & Pearson 相关系数（池化与日聚合两个尺度）
      & 池化：与负载 \(+0.508\)、净负荷 \(+0.260\)、光伏 \(+0.028\)；日聚合：与负载 \(+0.898\)、
        光伏 \(-0.343\)
      & \textbf{决定必须生成电价与净需求的联合情景}，否则系统性低估“高价时段同时缺电”的尾部风险 \\
    \bottomrule
  \end{tabularx}
\end{table}
\subsection{单参数灵敏度与 Sobol 全局灵敏度的完整档位}\label{sec:sens-app}

正文第九章给出各参数的扰动档位与“结论是否改变”的判断；本小节补充 Sobol 全局灵敏度的完整
指数，并说明它为什么必须\textbf{区分输出量}。

\begin{table}[htbp]
  \centering
  \caption{Sobol 全局灵敏度指数（典型日模型，$N=256$；输出量 1：典型日全天购电费，
           均值 $34\,603.42$ 元；输出量 2：相对无储能基线的节省率，均值 $27.02\%$；
           指数为量纲一）}
  \label{tab:sobol-app}
  \footnotesize
  \setlength{\tabcolsep}{5pt}
  \begin{tabular}{lcc}
    \toprule
    参数 & 全天购电费的一阶指数 \(S_1\) & 节省率的一阶指数 \(S_1\) \\
    \midrule
    储电量上限 \(\overline S\)        & 0.2510 & \(\le0.0033\) \\
    储电量下限 \(\underline S\)      & 0.1099 & \(\le0.0033\) \\
    往返/单边效率 \(\eta_c,\eta_d\)  & 0.0695 & \(\le0.0033\) \\
    充放功率上限 \(\overline P\)     & 0.0263 & \(\le0.0033\) \\
    电价水平                         & 0 & 0.8464 \\
    负载水平                         & 0 & 0.4711 \\
    光伏出力水平                     & 0 & \(\le0.0033\) \\
    峰谷价差                         & 0 & \(\le0.0033\) \\
    \bottomrule
  \end{tabular}
  \\[1ex]
  {\footnotesize 读法：对\textbf{全天购电费}起主导作用的是储电量上下限与效率；
   对\textbf{节省率}起主导作用的却是电价水平与负载水平——在问题 1 的确定性模型中，
   电价、光伏与负载水平对前者只是整体缩放，不改变调度结构，故一阶指数为 0。
   结论是“哪一个参数决定哪一个输出量”，而非“所有参数都重要”。}
\end{table}

\subsection{容量与充放功率的双向扫描明细}\label{sec:scan-app}

正文第四章给出扫描结论（容量摆幅是主要瓶颈）；本小节给出完整档位与两个单自由度放开实验。

\begin{table}[htbp]
  \centering
  \caption{储能容量与充放功率上限的双向扫描（问题 1，典型日；费用单位 元；
           “相对主口径”以主口径最优费用 $35\,126.95$ 元为分母，正号表示更省；
           容量扫描同步移动储电量上下限以保持 10\%--90\% 放电深度）}
  \label{tab:scan-app}
  \footnotesize
  \setlength{\tabcolsep}{5pt}
  \begin{tabular}{rrr|rrr}
    \toprule
    容量（kWh） & 最优购电费 & 相对主口径 & 充放功率上限（kW） & 最优购电费 & 相对主口径 \\
    \midrule
    7\,200  & 38\,640.45 & \(-10.00\%\) & 1\,250  & 40\,166.17 & \(-14.35\%\) \\
    9\,600  & 36\,700.45 & \(-4.48\%\)  & 2\,500  & 36\,390.84 & \(-3.60\%\) \\
    12\,000 & \textbf{35\,126.95} & \(0.00\%\) & 5\,000 & \textbf{35\,126.95} & \(0.00\%\) \\
    14\,400 & 33\,875.57 & \(+3.56\%\)  & 7\,500  & 35\,027.13 & \(+0.28\%\) \\
    19\,200 & 31\,692.43 & \(+9.78\%\)  & 10\,000 & 34\,960.59 & \(+0.47\%\) \\
    28\,800 & 28\,854.53 & \(+17.86\%\) & ---     & ---        & --- \\
    48\,000 & 27\,601.80 & \(+21.42\%\) & ---     & ---        & --- \\
    \bottomrule
  \end{tabular}
  \\[1ex]
  {\footnotesize 另做两个单自由度放开实验：仅解除储电量上限（充放功率维持 $5000$~kW）为
   $27\,351.55$~元，相对主口径多省 $22.14\%$；仅解除充放功率上限（容量维持 $12\,000$~kWh）
   为 $34\,544.43$~元，多省 $1.66\%$。界活跃统计：充电功率界在 $19/144$ 个时段被打满、
   放电功率界 $0/144$，储电量上界活跃 $21/144$、下界活跃 $8/144$。}
\end{table}

\subsection{L 形（Benders）分解与直接求解的独立对拍}\label{sec:lshape}

问题 2--4 的模型都是\textbf{两阶段结构}：第一阶段是全场共享的计划量 $(E_t,c_t,d_t,S_t)$，
第二阶段是逐情景的追索量 $g_{\omega,t}$。本文主结果全部用 HiGHS 直接求解（单日 $4n+Kn$ 个变量，
$K=60$ 时 $9216$ 个）；但“求解器返回最优”本身不构成模型正确性的证据，因此本节\textbf{另写一套
分解算法}作为对照：把问题改写为 L 形（Benders）形式，用与直接求解\textbf{完全不同的数值路径}
再算一遍。给定第一阶段 $(E,c,d,S)$ 后，情景 $\omega$ 的追索问题逐槽完全可分：
\begin{equation}
  Q_\omega(E,c,d)=\sum_t 5p_t\max\bigl\{0,\ \widetilde{ND}_{\omega,t}+c_t-d_t-E_t\bigr\},
  \label{eq:lshape-sub}
\end{equation}
每个槽的最优值就是一个正部函数，对 $E_t$ 的次梯度为 $-5p_t$（缺口为正时）或 $0$（否则），
故\textbf{每个情景每次迭代可解析地给出 144 条最优性割}（多割），无需调用求解器求子问题；
主问题只保留 $4n+1=577$ 个变量。

\begin{table}[htbp]
  \centering
  \caption{L 形分解与直接求解的对照（对象：单日两阶段随机规划，题面指定四个代表日
           （2025-03-20、06-21、09-23、12-21）的结果合计；范围：$s_{init}=6000$~kWh、
           终端模式 $S_n\ge s_{init}$、两法共用同一份分层 medoid 情景；
           时间列单位为 s（四日合计），变量数单位为个；
           相对间隙 $=|z_{\text{L 形}}-z_{\text{直接}}|/|z_{\text{直接}}|$；
           数据源：\texttt{intermediate/decomp.json} 的 \texttt{scenario\_scan} 字段）}
  \label{tab:lshape}
  \small
  \setlength{\tabcolsep}{4pt}
  \begin{tabular}{crrrrrc}
    \toprule
    情景数 $K$ & 直接求解变量 & L 形变量 & 直接求解 & L 形分解 & 加速比 & 最大相对间隙 \\
    \midrule
    $20$  & $3\,456$ & $720$ & $0.142$ & $0.281$ & $0.50\times$ & $6.65\times10^{-16}$ \\
    $60$  & $9\,216$ & $720$ & $0.356$ & $0.358$ & $1.00\times$ & $2.09\times10^{-15}$ \\
    $120$ & $17\,856$ & $720$ & $1.291$ & $0.401$ & $3.22\times$ & $3.06\times10^{-15}$ \\
    $300$ & $43\,776$ & $720$ & $3.721$ & $0.635$ & $5.86\times$ & $2.60\times10^{-15}$ \\
    $500$ & $72\,576$ & $720$ & $8.962$ & $0.875$ & $10.25\times$ & $4.60\times10^{-15}$ \\
    \bottomrule
  \end{tabular}
\end{table}

\textbf{三条结论。} ① \textbf{两条数值路径给出的最优值一致到机器精度}：结果窗口全部 334 天上
（$K=60$）逐日对比最大相对间隙 $3.76\times10^{-15}$、全部收敛（平均迭代 9.75 次、
单日最多 1097 条割），$K=300$ 的 42 天子样本上为 $4.82\times10^{-15}$——两个独立实现同时给出
同一最优值，比“求解器说 optimal”强得多。② \textbf{加速比随情景数增长}：$K=20$ 时分解反而慢约
一倍，$K\approx60$ 附近是分界点，$K=120$、300、500 时分别加速 3.22、5.86、10.25 倍；
主问题变量数始终为 720、不随 $K$ 增长，这是加速的来源。③ \textbf{工程取舍}：本题主口径 $K=60$、
单日直接求解 0.106~s（L 形 0.092~s，两者同量级），分解法在该规模上\textbf{没有}速度优势；
把解析割聚合成单割后 200 次迭代仍未收敛、相对间隙 0.65\%--3.42\%，而多割 9--14 次即收敛，
故\textbf{主结果仍用直接求解}；分解法的价值在于给出独立最优性对拍，并在情景数放大到数百时
把单日求解稳定在 0.8~s 以内。\evid{C-Q5-01}







\begin{table}[htbp]
  \centering
  \caption{光伏预报偏差的逐整点结构（样本为 2025 年每天 \(0{:}00\) 发布的 24~h 预报与
           附件 2 对应实际功率的逐槽配对，\(n=52\,560\)；bias 单位 kW，
           “极值整点”给出该组内绝对偏差最大的整点及其取值）}
  \label{tab:biasprofile}
  \footnotesize
  \setlength{\tabcolsep}{4pt}
  \begin{tabular}{lccc}
    \toprule
    整点分组 & 整点范围 & 组内 bias 区间 & 极值整点 \\
    \midrule
    夜间（光伏无有效出力） & \(0\)--\(4\) 时 & \(0.00\) 至 \(+0.09\) & \(4\) 时 \(+0.09\) \\
    上午（系统性\textbf{高估}） & \(5\)--\(10\) 时 & \(+214.21\) 至 \(+733.00\) & \(7\) 时 \(+733.00\) \\
    正午过渡 & \(11\) 时 & \(+7.21\) & --- \\
    下午（系统性\textbf{低估}） & \(12\)--\(18\) 时 & \(-801.87\) 至 \(-106.46\) & \(15\) 时 \(-801.87\) \\
    傍晚之后（预报空洞） & \(19\)--\(23\) 时 & \(-0.15\) 至 \(0.00\) & \(19\) 时 \(-0.15\) \\
    \bottomrule
  \end{tabular}
\end{table}





\section{场景保真度与电价检验明细}

本节给出正文第五章场景集与第七章电价过程检验、CVaR 风险—费用前沿的完整明细表。

（联合情景保真度检验表见 \ref{tab:scenfid}。）



\begin{table}[htbp]
  \centering
  \caption{风险前沿三个代表点的情景内指标与真实轨迹回测
           （对象：联合情景集上按 \(\lambda\) 求解的三组解；范围：2025-02-01 至 12-31 中每 6 天取 1 天的 56 个代表日；
            单位：费用 元/日、电量 kWh/日；样本：64 个代表情景，\(\alpha=0.95\)）}
  \label{tab:frontier-backtest}
  \footnotesize
  \setlength{\tabcolsep}{4pt}
  \begin{tabularx}{\linewidth}{L{3.0cm} C{2.3cm} C{2.3cm} C{2.5cm} Y}
    \toprule
    代表点 & 情景内期望费用 & 情景内 95\% CVaR & 真实轨迹回测费用 & 紧急购电量 \\
    \midrule
    \(\lambda=0\)（期望最优） & 41\,995.59 & 52\,176.36 & 45\,694.27 & 480.19 \\
    \(\lambda=0.3\)（推荐） & 42\,168.43 & 50\,814.83 & 45\,778.03 & 324.60 \\
    \(\lambda=1\)（近似最坏） & 42\,411.56 & 50\,510.46 & 45\,952.85 & 313.83 \\
    \bottomrule
  \end{tabularx}
\end{table}

\section{验证映射、风险结果与局限明细}

本节给出正文第九章验证层级映射、第八章问题 4 逐日期结果与第十章局限清单的完整表格。





\section{附件结构与问题 4 逐日期结果明细}

本节给出四个附件的字段结构与问题 4 在四个指定日期的逐项结果。



【问题 4 逐日期结果】




\section{明细表与口径框（正文摘要引用）}
本节集中给出正文以摘要形式引用的明细表与全部核心指标口径框（\calref{C-01}--\calref{C-11}），便于评审核对；正文对应处均给出结论、口径编号与本节的表号。













\begin{caliber}{C-02}{全天购电量}
\textbf{定义}\quad \(E^{\text{day}}=\sum_{t=1}^{144}E_t\)，即该日 144 个 10~min 槽的\textbf{计划购电量}之和；%
\textbf{公式}\quad 结果为逐槽求和，无闭式；结果文件每行末列即该日汇总；%
\textbf{数据}\quad 结果文件「计划购电量」工作表的 144 个时段列；%
\textbf{时段}\quad 单日；问题 1 为典型日，问题 2--4 为 2025-02-01 至 12-31 的 334 天；%
\textbf{单位}\quad kWh（结果文件保留 4 位小数，正文保留 2--4 位）；%
\textbf{含}\quad 问题 2、4-2 中即第一阶段计划量；问题 3、4-3 中调整后的量一律称“全天调整购电量”并单独命名；%
\textbf{不含}\quad 不含储能充电量、不含紧急购电量 \calref{C-05}、不含盈余处置量 \calref{C-06}；%
\textbf{对照}\quad 问题 2、3 的计划购电量模型内值为 \(22\,220\,333.2033\)~kWh，结果文件回读值为 \(22\,220\,333.1391\)~kWh，
相对差 \(2.9\times10^{-9}\)，来自写入时的 4 位小数取整；%
\textbf{复算}\quad \texttt{results\_io.py::plan\_sheet\_frame()}（写出行汇总列）与 \texttt{audit.py::audit\_result\_files()}（回读结果文件后重算 \(\sum_{144}\) 槽）
\end{caliber}

\begin{caliber}{C-03}{总购电费用（计划 \(+\) 紧急 \(+\) 调整相关）}
\textbf{定义}\quad \(C^{\text{total}}=C^{\text{plan}}+C^{\text{breach}}+C^{\text{excess}}+C^{\text{emg}}\)；
问题 2、4-2 无调整机制，后两项恒为 0；%
\textbf{公式}\quad 问题 2、4-2：\(C=\sum_tp_tE_t+5\sum_tp_tg_t\)；
问题 3、4-3：\(C^{\text{plan}},C^{\text{breach}},C^{\text{excess}}\) 按 \calref{C-09} 的偏差量定价口径计算，
\(C^{\text{emg}}=\sum_t5p_tg_t\)；%
\textbf{数据}\quad 电价取附件 1（问题 1--3）或附件 4（问题 4）；电量取本队结果文件；%
\textbf{时段}\quad 问题 1 为典型日；问题 2--4 为 334 天；%
\textbf{单位}\quad 元（保留 2 位小数，内部计算不取整）；%
\textbf{不含}\quad 不含储能投资与运维成本、容量补偿、辅助服务收益、绿电环境价值、循环衰减与安全成本；%
\textbf{对照}\quad 若只报计划购电费而不计紧急购电，问题 2 为 \(13\,637\,552.90\)~元，
比 \(C^{\text{total}}\) 低 \(627\,239.24\)~元（\(4.40\%\)）；%
\textbf{复算}\quad \texttt{q2.py::run\_q2()}、\texttt{q3.py::summary()}、\texttt{q4.py::run\_q4\_2()} 与 \texttt{audit.py::audit\_result\_files()}，
四项之和与总费用的偏差为 0
\end{caliber}

\begin{caliber}{C-04}{节省率}
\textbf{定义}\quad \(\eta=\bigl(C^{base}-C^{\text{total}}\bigr)/C^{base}\times100\%\)，分母与分子必须同窗口、同电价口径；%
\textbf{公式}\quad \(C^{base}\) 取 \calref{C-01} 的对应变体（V1 典型日 / V2 固定价窗口 / V3 波动价窗口，
变体定义见正文表~\bref{tab:cbases}）；%
\textbf{数据}\quad \(C^{base}\) 与 \(C^{\text{total}}\) 均取自本队结果；%
\textbf{时段}\quad 与 \calref{C-01} 变体一致；%
\textbf{单位}\quad \%（保留 2 位小数）；%
\textbf{不含}\quad 不得相对贪心策略或完全信息下界报同一个数而不加限定；相对其他基线的结果必须\textbf{单独命名}；%
\textbf{对照}\quad 问题 1 相对 V3 为 \(26.90\%\)；问题 2、3 相对 V1 为 \(13.06\%\)、\(15.08\%\)；
问题 4-2、4-3 相对 V2 为 \(11.35\%\)、\(13.54\%\)；%
\textbf{复算}\quad \texttt{paper\_numbers.py::main()}（由 \calref{C-01} 与 \calref{C-03} 导出）
\end{caliber}

\begin{caliber}{C-06}{弃光/盈余处置电量}
\textbf{定义}\quad 光伏盈余与超购部分的处置量：典型日按
\(s_t=PV_t\Delta t+E_t+d_t-L_t\Delta t-c_t\ge0\) 逐槽统计（式~\bref{eq:q1-balance} 的松弛量）；
窗口口径按实际数据复算为 \(\sum_t\max\{PV_t\Delta t+E_t+d_t+g_t-L_t\Delta t-c_t,0\}\)；%
\textbf{公式}\quad 上式为定义式；该量非负，其余额为超购部分被弃置的电量；%
\textbf{数据}\quad \(L,PV\) 取附件 2 实际值，其余取本队结果文件；%
\textbf{时段}\quad 逐槽统计、按窗口（334 天）汇总；%
\textbf{单位}\quad kWh（按电量统计，\textbf{不得}按功率报）；%
\textbf{不含}\quad 不含“因充电功率受限而未发生的充电”，不含未购电部分；%
\textbf{对照}\quad 问题 1 典型日最优弃光 \(=0.0000\)~kWh（数值残差 \(1.37\times10^{-12}\)~kWh）；
问题 2、3 的窗口盈余处置量为 \(3\,106\,390.84\)~kWh 与 \(2\,638\,963.35\)~kWh；
问题 4-2、4-3 的同一量为 \(3\,100\,463.79\)~kWh 与 \(2\,645\,427.81\)~kWh（均见正文表~\bref{tab:summary}）；%
\textbf{复算}\quad \texttt{config.py::curtailment(price, load, pv, E, c, d)}（典型日）与 \texttt{audit.py::audit\_internal()}（窗口口径）
\end{caliber}

\begin{caliber}{C-07}{充放电量与储电量（交流侧）}
\textbf{定义}\quad 结果文件中的 \(c_t,d_t\) 为\textbf{并网点（AC 侧）}电量；储电量为\textbf{时段末}值 \(S_t\)；%
\textbf{公式}\quad 递推见式~\bref{eq:q1-soc}，\(\eta_c=\eta_d=0.9\)（单边口径，往返 \(0.81\)）；
充放功率上限 \(5000\)~kW 对应单槽上限 \(\overline P\Delta t=833.3333\)~kWh；%
\textbf{数据}\quad 本队结果文件「充放电量」工作表；%
\textbf{时段}\quad 逐槽；表 2 格式按 6 个 4~h 块（每块 24 槽）汇总，并给出 \(0{:}00\) 与 \(24{:}00\) 两个储电量点；%
\textbf{单位}\quad kWh；运行区间 \([1200,10800]\)~kWh；%
\textbf{不含}\quad 不含电池侧（DC）电量、不含自放电与循环衰减（假设 A4）；%
\textbf{对照}\quad 恒等式校验：问题 1 的充放比 \(20\,740.67/16\,799.94=1.2346=1/(\eta_c\eta_d)\)，
与 \(\eta_c\eta_d=0.81\) 互为倒数，偏差 \(4.44\times10^{-16}\)；问题 2 四个指定日期的合计充放比同为 \(1.2346\)；%
\textbf{复算}\quad \texttt{results\_io.py::soc\_sheet\_frame()}（块级汇总）与 \texttt{audit.py::audit\_result\_files()}（块级净递推）
\end{caliber}

\begin{caliber}{C-08}{储电量递推与终端条件}
\textbf{定义}\quad 储电量\textbf{跨日连续递推}（\(S_{d,0}=S_{d-1,144}\)，假设 A5），
窗口末端施加终端条件 \(S_T\ge S_0\)（假设 A6），问题 1 另加严格日周期 \(S_0=S_{144}\)；%
\textbf{公式}\quad \(S_t=S_{t-1}+\eta_cc_t-d_t/\eta_d\)，\(S_t\in[1200,10800]\)~kWh；
2025 年 1 月为预热期（假设 A8），其终值即 2 月 1 日初值；%
\textbf{数据}\quad 结果文件「充放电量」工作表的储电量列与中间结果 \texttt{q2\_main.json} 等；%
\textbf{时段}\quad 逐槽；窗口 2025-02-01 至 12-31（334 天）；%
\textbf{单位}\quad kWh；%
\textbf{不含}\quad 不含“每日清零”“每日重置为 6000”两种读法——前者会引入末日掏空型短视解，后者与 1 月预热期矛盾；%
\textbf{对照}\quad 问题 1：日周期口径 \(35\,126.9486\)~元（\(S_0=S_{144}=6000.00\)~kWh）；
仅要求两端相等的自由环状口径 \(35\,118.5986\)~元（最优 \(S=8550.00\)~kWh），相差 \(8.35\)~元（\(0.02\%\)）；
窗口口径下“末值低于初值的天数”为 0；%
\textbf{复算}\quad \texttt{lp\_engine.py::solve\_day()}（终端约束）与 \texttt{audit.py::audit\_result\_files()}（\texttt{A6\_末值低于初值天数}）
\end{caliber}



\section{问题 2--4 的逐时段结果明细表}
第八章已按题面表 1／表 2／表 3 的格式给出四个指定日期的结果；本节补充\textbf{指定 10~min 时段}的购电量明细与逐时段充放电明细，便于与结果文件逐格对照。

\begin{table}[htbp]
  \centering
  \caption{微网在指定时段与指定日期的购电量、全天购电量和购电费（问题 2；电量单位：\unit{kWh}，费用单位：\unit{元}）}
  \label{tab:q2-buy}
  \small
  \begin{tabular}{lrrrr}
    \toprule
    时段／指标 & 2025-03-20 & 2025-06-21 & 2025-09-23 & 2025-12-21 \\
    \midrule
    \(10{:}00\)--\(10{:}10\) & 0.0000 & 0.0000 & 0.0000 & 0.0000 \\
    \(12{:}00\)--\(12{:}10\) & 804.3029 & 0.0000 & 536.2415 & 997.0335 \\
    \(14{:}00\)--\(14{:}10\) & 0.0000 & 0.0000 & 0.0000 & 0.0000 \\
    \(16{:}00\)--\(16{:}10\) & 817.3061 & 255.2254 & 902.7969 & 1001.2240 \\
    \(18{:}00\)--\(18{:}10\) & 697.9973 & 0.0000 & 764.3298 & 715.4144 \\
    \(20{:}00\)--\(20{:}10\) & 0.0000 & 0.0000 & 0.0000 & 0.0000 \\
    \midrule
    全天购电量 & 74\,178.1298 & 35\,701.5391 & 74\,744.7963 & 99\,692.6601 \\
    全天购电费 & 45\,229.90 & 20\,324.26 & 47\,063.39 & 64\,376.98 \\
    紧急购电量 & 1025.8805 & 187.0843 & 193.3669 & 22.5250 \\
    全天总费用 & 49\,654.02 & 21\,255.45 & 47\,717.70 & 64\,432.28 \\
    \bottomrule
  \end{tabular}
  \\[1ex]
  {\footnotesize 表注：时段按题面表述取区间左端时刻；
  全天购电费为计划口径（\(\sum_t p_tE_t\)），全天总费用为计划购电费与紧急购电费之和；
  四个日期的无储能基线分别为 52\,290.65、29\,668.18、54\,741.30、71\,176.92~\unit{元}。}
\end{table}

\begin{table}[htbp]
  \centering
  \caption{储能设备在指定日期的充放电量及 \(0{:}00\) 与 \(24{:}00\) 的储电量（问题 2，单位：\unit{kWh}）}
  \label{tab:q2-soc}
  \footnotesize
  \begin{tabular}{L{2.0cm}L{2.1cm}rr@{\quad}L{2.0cm}L{2.1cm}rr}
    \toprule
    日期 & 时段 & 充电量 & 放电量 & 日期 & 时段 & 充电量 & 放电量 \\
    \midrule
    2025-03-20 & \(0{:}00\)--\(4{:}00\)   & 4\,500.0000 & 0.0000    & 2025-09-23 & \(0{:}00\)--\(4{:}00\)   & 4\,500.0000 & 0.0000 \\
    2025-03-20 & \(4{:}00\)--\(8{:}00\)   & 833.3333   & 6\,224.1108 & 2025-09-23 & \(4{:}00\)--\(8{:}00\)   & 833.3333   & 5\,547.8264 \\
    2025-03-20 & \(8{:}00\)--\(12{:}00\)  & 6\,500.0000 & 2\,415.8892 & 2025-09-23 & \(8{:}00\)--\(12{:}00\)  & 5\,833.3333 & 3\,020.9227 \\
    2025-03-20 & \(12{:}00\)--\(16{:}00\) & 6\,905.5931 & 2\,218.5304 & 2025-09-23 & \(12{:}00\)--\(16{:}00\) & 6\,728.4081 & 1\,606.2615 \\
    2025-03-20 & \(16{:}00\)--\(20{:}00\) & 0.0000    & 5\,709.7835 & 2025-09-23 & \(16{:}00\)--\(20{:}00\) & 0.0000    & 5\,602.9091 \\
    2025-03-20 & \(20{:}00\)--\(24{:}00\) & 5\,333.3333 & 2\,930.2165 & 2025-09-23 & \(20{:}00\)--\(24{:}00\) & 5\,333.3333 & 3\,037.0909 \\
    \midrule
    2025-06-21 & \(0{:}00\)--\(4{:}00\)   & 0.0000    & 0.0000    & 2025-12-21 & \(0{:}00\)--\(4{:}00\)   & 4\,500.0000 & 0.0000 \\
    2025-06-21 & \(4{:}00\)--\(8{:}00\)   & 140.7851  & 4\,303.4721 & 2025-12-21 & \(4{:}00\)--\(8{:}00\)   & 1\,666.6667 & 1\,508.3333 \\
    2025-06-21 & \(8{:}00\)--\(12{:}00\)  & 5\,207.8672 & 130.5638 & 2025-12-21 & \(8{:}00\)--\(12{:}00\)  & 5\,833.3333 & 7\,806.6667 \\
    2025-06-21 & \(12{:}00\)--\(16{:}00\) & 5\,458.7995 & 0.0000    & 2025-12-21 & \(12{:}00\)--\(16{:}00\) & 8\,333.3333 & 2\,835.0000 \\
    2025-06-21 & \(16{:}00\)--\(20{:}00\) & 0.0000    & 6\,154.1006 & 2025-12-21 & \(16{:}00\)--\(20{:}00\) & 0.0000    & 5\,797.4742 \\
    2025-06-21 & \(20{:}00\)--\(24{:}00\) & 5\,333.3333 & 2\,485.8994 & 2025-12-21 & \(20{:}00\)--\(24{:}00\) & 5\,333.3333 & 2\,842.5258 \\
    \bottomrule
  \end{tabular}
  \\[1ex]
  {\footnotesize 表注：四个指定日期的 \(0{:}00\) 与 \(24{:}00\) 储电量均为 6000.0000~\unit{kWh}；
  每段含 24 个 10~\unit{min} 槽，充放电量为交流侧电量（\calref{C-07}）；
  四个日期的放电量／充电量之比均等于往返效率 \(\eta_c\eta_d=0.8100\)，
  其倒数即充电量／放电量之比 \(1/(\eta_c\eta_d)=1.2346\)。
  本表与正文表~\bref{tab:q2-t2} 取自同一数据源 \texttt{results/result2.xlsx}，逐格一致。}
\end{table}

\begin{table}[htbp]
  \centering
  \caption{微网在指定日期的紧急购电量（问题 2；购电量单位：\unit{kWh}，电价单位：\unit{元/kWh}）}
  \label{tab:q2-emergency}
  \footnotesize
  \begin{tabular}{llrr}
    \toprule
    日期 & 购电时间段 & 紧急购电量 & 交易时刻电价（加权平均） \\
    \midrule
    2025-03-20 & \(6{:}00\)--\(8{:}00\)   & 228.9914 & 0.8625 \\
    2025-03-20 & \(8{:}10\)--\(8{:}30\)   & 49.1586  & 0.8625 \\
    2025-03-20 & \(8{:}50\)--\(10{:}30\)  & 531.2739 & 0.8625 \\
    2025-03-20 & \(10{:}40\)--\(11{:}40\) & 216.4567 & 0.8625 \\
    \midrule
    2025-06-21 & \(4{:}00\)--\(4{:}50\)   & 5.2416   & 0.9955 \\
    2025-06-21 & \(14{:}00\)--\(14{:}30\) & 181.8427 & 0.9955 \\
    \midrule
    2025-09-23 & \(6{:}50\)--\(7{:}40\)   & 132.5948 & 0.6768 \\
    2025-09-23 & \(11{:}00\)--\(11{:}40\) & 60.7721  & 0.6768 \\
    \midrule
    2025-12-21 & \(11{:}00\)--\(11{:}30\) & 22.5250 & 0.4910 \\
    \midrule
    \multicolumn{2}{l}{四个日期合计} & 1428.8567 & --- \\
    \bottomrule
  \end{tabular}
  \\[1ex]
  {\footnotesize 表注：本表按“日期 \(\times\) 紧急购电事件”逐行展开，与题面表 3 的横向版式一一对应；
  同一日期可有多行，连续触发的相邻槽按事件聚合，时间段为该事件涉及的槽位区间；
  交易时刻电价为该事件内按电量加权的平均电价（事件电价按 \(5p_t\) 计价，本列为 \(p_t\)）。
  无紧急购电的时段不列入本表（全窗口 334 天中 296 天有触发）。}
\end{table}

\section{口径声明与时段对照明细表}




\section{伪验证自查清单（逐条证据）}\label{sec:fakeverify-app}

\noindent
\textbf{① 自证循环}：审计脚本不共用建模费用函数，费用双路径复算偏差 \(4.32\times10^{-7}\)；\\
\textbf{② 以求解器状态代替物理校验}：求解器返回 optimal 只作必要条件，同时报告供给缺口、储电量递推、
越界与互斥等物理残差；\\
\textbf{③ 比较不同样本}：R1／R2／R3、\(n=1/2/4\)、\(K=60/120\)、四种情景生成器均在同一场景集或同一
实际轨迹上配对比较；\\
\textbf{④ 训练集当测试集}：主结果明确声明为样本内评估，另做样本外回测（\bref{sec:oos}）与生成器样本外对照（\ref{sec:q4-gen}）；\\
\textbf{⑤ 只报点估计}：关键量均给多档对照——效率 4 档、容量 7 档、功率 5 档、\(\kappa\) 5 档、
场景数 2 档、CVaR 的 \(\lambda\) 11 档、机会约束 \(\beta\) 6 档、盒式鲁棒 \(k\) 7 档；\\
\textbf{⑥ 事后挑选指标}：本节指标在建模前由口径 C-01 至 C-11 固定，未随结果调整；\\
\textbf{⑦ 把“结果好看”当“模型正确”}：问题 2 的 \(13.06\%\) 与问题 1 的 \(26.90\%\) 同量级且更小，
差额可由预报误差成本 \(853\,231.23\)~元解释，未把“节省率更高”当作模型更好的证据；\\
\textbf{⑧ 扰动无效}：灵敏度扰动作用于运行区间与功率上限本身，且每档重新求解，不复用已有解。

\section{附图}

\begin{figure}[htbp]
  \centering
  \includefig[width=0.86\linewidth]{fig12_scenario_generators.pdf}
  \caption{四种情景生成器的样本外对照（左：相对“历史代表日”的总费用差；
           中：184 天累计紧急购电量；右：日总残差的 1\%/99\% 分位最大绝对偏差；
           对象：训练窗 2025-02-01 至 06-30、测试窗 07-01 至 12-31 共 184 天；
           单位：百分比、kWh；样本：情景数由二维分层固定为 48 个，Copula 生成器抽样 600 条后削减）}
  \label{fig:gen-cmp}
\end{figure}
